\documentclass[12pt,a4paper]{article}
\usepackage[utf8]{inputenc}
\usepackage[T1]{fontenc}
\usepackage{amsmath, amssymb, mathtools}
\usepackage{graphicx}
\usepackage{xcolor}
\usepackage{geometry}
\usepackage[colorlinks=true, allcolors=blue]{hyperref}
\usepackage{amsthm}
\newtheorem{theorem}{Theorem}[section]
% \newtheorem{corollary}{Corollary}[corollary]
\newtheorem{lemma}[theorem]{Lemma}
\newtheorem{corollary}[theorem]{Corollary}
\newtheorem{definition}[theorem]{definition}
\geometry{
    top=2cm,
    bottom=3cm,
    left=1cm,
    right=1cm
}
\usepackage{float}
\usepackage{cancel}
\usepackage{esint}
\usepackage{mathdots}
\usepackage{mhchem}
\usepackage{stackrel}
\usepackage{stmaryrd}
% Define custom colors if needed (from original preamble)
\definecolor{blue}{RGB}{14,107,217}
\definecolor{green}{RGB}{0,158,40}
\definecolor{red}{RGB}{235,16,16}
\definecolor{brown}{RGB}{164,66,0}
\definecolor{orange}{RGB}{231,135,26}
\definecolor{purple}{RGB}{94,53,177}
% Custom macros from the original LyX file
\newcommand{\R}{\mathbb{R}}
\newcommand{\N}{\mathbb{N}}
\newcommand{\BB}{\mathcal{B}}
\newcommand{\Q}{\mathbb{Q}}
\newcommand{\Z}{\mathbb{Z}}
\newcommand{\C}{\mathbb{C}}
\renewcommand{\O}{\Omega}
\newcommand{\eps}{\varepsilon}
\begin{document}
\begin{center}
    \fbox{\parbox{\dimexpr\textwidth-2\fboxsep-2\fboxrule\relax}{\centering \Large \textbf{Calculate M}}}
\end{center}
\tableofcontents 
\newpage
\vspace{0.5cm}
\section{Problem Definition}

Let $\mathcal{A}\in \R^{d\times d \times \cdots \times d}$  be an $MPOP$,  with $N$ vertices of order $k$ and a physical dimension $d$. We will calculate the error of the sketching tensor network $S$, that is: 
\begin{align}
\text{var}(\text{Tr}(S_{1}AS_{2}^{T}S_{2}A\cdots S_{k}^{T}S_{k}AS_{1}^{T})) & =\mathbb{E}[\text{Tr}(S_{1}A\cdots S_{k}^{T}S_{k}AS_{1}^{T})^{2}]-\mathbb{E}[\text{Tr}(S_{1}A\cdots S_{k}^{T}S_{k}AS_{1}^{T})]^{2}\\
 & =\mathbb{E}[\text{Tr}(S_{1}A\cdots S_{k}^{T}S_{k}AS_{1}^{T})^{2}]-\text{Tr}(A^{k})^{2}.\label{eq:target_eq}
\end{align}
We evaluate the 1st term: 
\begin{align}
\mathbb{E}[\text{Tr}(S_{1}AS_{2}^{T}\cdots S_{k}^{T}S_{k}AS_{1}^{T})^{2}] & =\mathbb{E}[\text{Tr}(S_{1}A\cdots S_{k}^{T}S_{k}AS_{1}^{T})^{2}]\nonumber\\
 & =\mathbb{E}[\text{Tr}(\prod_{i=1}^{k}S_{i}^{T}S_{i}A)^{2}]\nonumber\\
 & =\mathbb{E}[\text{Tr}(\prod_{i=1}^{k}S_{i}^{T}S_{i}A\otimes\prod_{i=1}^{k}S_{i}^{T}S_{i}A)]\nonumber\\
 & =E\left[\text{Tr}\left(\prod_{i=1}^{k}\left(S_{i}^{T}S_{i}\otimes S_{i}^{T}S_{i}\right)\cdot A\otimes A\right)\right]\label{eq:k_target}
 \end{align}
 

where the last transition holds since the $S_{i}$ are $i.i.d.$
We denote 
\begin{align}
       M_{i^{(a)}i^{(b)}i^{(c)}i^{(d)}}&:=\mathbb{E}[(S^{T}S)_{i^{(a)}i^{(b)}}\otimes(S^{T}S)_{i^{(c)}i^{(d)}}]\\
       B&:=A \otimes A \label{eq:M_def}
\end{align}
We note that $M,B\in \mathbb{R}^{d^{2N}\times d^{2N}}$


and we wish to evaluate 
\[
Tr(M \cdot B)
\]
and 
\[
Tr((MB)^k).
\]

\newpage

\subsection{The Structure of sketching Tensor Network (TN) $S$}

The tensor $S\in\R^{m\times d\times \cdots \times d}$ of order $N+1$ has the following form:

\vspace{0.5cm}
\begin{center}
\begin{figure}[H]
    \centering
    \includegraphics[width=0.5\linewidth]{sketch_of_graph_network.png}
    \caption{Graphical representation of $S$}
    \label{fig:s_structure}
\end{figure}
\end{center}
\vspace{0.5cm}

In the current graph $X^{(i)}\in\R^{m\times d}$ there are $i.i.d.$ matrices with $X_{ij}\sim \mathcal{N}(0,\sigma^{2})$ that represent the Kronecker part of the sketching. Let $g^{(i)}\in\R^{m\times m_{1}\times m_{2}\times m}$ are the contractions of the original gadgets that we reshape to an order 3 tensor with $g^{(i)}\in\R^{m\times m\times m}$ where $m_{1}\cdot m_{2}=m$. 
Our gadgets $g^{(i)}$  are constructed with the following:
\begin{align}
G_{ij} & \sim \mathcal{N}(0,\sigma^{2})\\
g_{k_{1}k_{2}^{(1)}k_{2}^{(2)}z_{1}}^{(1)} & =\sum_{l=1}^{m}G^{(1,1)}_{k_{1}k_{2}^{(1)}l}\cdot G^{(1,2)}_{lk_{2}^{(2)}z_{1}}\\
g_{z_{p-1}k_{p+1}^{(1)}k_{p+1}^{(2)}z_{p}}^{(p)} & =\sum_{l=1}^{m}G^{(p,1)}_{z_{p-1}k_{p+1}^{(1)}l}G^{(p,2)}_{lk_{p+1}^{(2)}z_{p}} \label{eq:g_def}
\end{align}

The matricization $S\in\R^{m\times d^{N}}$, for all $1\le n\le m, 1\le i\le d^{N}$, where $i=(i_{1},...,i_{N})\in[d]^{N}$ can be written as:


\begin{align}
S_{n,(i_{1},...,i_{N})} & =\sum_{z_{1},.,z_{N-2},k_{1},.,k_{N}=1}^{m}g_{k_{1}k_{2}z_{1}}^{(1)}g_{z_{1}k_{3}z_{2}}^{(2)}\cdots g_{z_{N-2}k_{N}n}^{(N-1)}\cdot X_{i_{1}k_{1}}^{(1)}X_{i_{2}k_{2}}^{(2)}\cdots X_{i_{N}k_{N}}^{(N)}\\
 & =\sum_{\boldsymbol{z},\boldsymbol{k}=1}^{m}g_{k_{1}k_{2}z_{1}}^{(1)}\prod_{p=2}^{N-2}g_{z_{p-1}k_{p+1}z_{p}}^{(p)}\cdot g_{z_{N-2}k_{N}n}^{(N-1)}\prod_{l=1}^{N}X_{i_{l}k_{l}}^{(l)} \label{eq:s_def}
\end{align}


where we denote $\boldsymbol{z}=(z_{1},...,z_{N-2}), \boldsymbol{k}=(k_{1},...,k_{N})$ with $1\le z_{i},k_{i}\le m$,

and the columns of $S\in\R^{m\times d^{N}}$:
\[
S=\begin{bmatrix}| &  & | &  & |\\
S_{1} & \cdots & S_{i} & \cdots & S_{d^{N}}\\
| &  & | &  & |
\end{bmatrix}.
\]

We obtain $S^{T}S\in\R^{d^{N}\times d^{N}}$, where each entry is defined by
\begin{align}
\left(S^{T}S\right)_{i^{(a)}i^{(b)}} & =\langle S_{i^{(a)}},S_{i^{(b)}}\rangle=\sum_{n=1}^{m}S_{ni^{(a)}}S_{ni^{(b)}}\label{eq:dot_prod_example}\\
 & =\sum_{n=1}^{m}\left(\sum_{\boldsymbol{z},\boldsymbol{k}=1}^{m}g_{k_{1}k_{2}z_{1}}^{(1)}g_{z_{1}k_{3}z_{2}}^{(2)}\cdots g_{z_{N-2}k_{N}\boldsymbol{n}}^{(N-1)}\prod_{l=1}^{N}X_{i_{l}^{(a)}k_{l}}\sum_{\boldsymbol{z}^{\prime},\boldsymbol{k}^{\prime}=1}^{m}g_{k_{1}^{\prime}k_{2}^{\prime}z_{1}^{\prime}}^{(1)}g_{z_{1}^{\prime}k_{3}^{\prime}z_{2}^{\prime}}^{(2)}\cdots g_{z_{N-2}^{\prime}k_{N}^{\prime}\boldsymbol{n}}^{(N-1)}\prod_{l=1}^{N}X_{i_{l}^{(b)}k_{l}^{\prime}}\right)\nonumber\\
 & =\sum_{n=1}^{m}\left(\sum_{\boldsymbol{z},\boldsymbol{k},\boldsymbol{z}^{\prime},\boldsymbol{k}^{\prime}=1}^{m}g_{k_{1}k_{2}z_{1}}^{(1)}g_{k_{1}^{\prime}k_{2}^{\prime}z_{1}^{\prime}}^{(1)}\prod_{p=2}^{N-2}g_{z_{p-1}k_{p+1}z_{p}}^{(p)}g_{z_{p-1}^{\prime}k_{p+1}^{\prime}z_{p}^{\prime}}^{(p)}\cdot g_{z_{N-2}k_{N}\boldsymbol{n}}^{(N-1)}g_{z_{N-2}^{\prime}k_{N}^{\prime}\boldsymbol{n}}^{(N-1)}\prod_{l=1}^{N}X_{i_{l}^{(a)}k_{l}}\prod_{l=1}^{N}X_{i_{l}^{(b)}k_{l}^{\prime}}\right)\nonumber \label{eq:S_dot_prod_def}
\end{align}


Therefore, to evaluate the Kronecker part of the sketch \ref{eq:s_def}, we get $4$ copies of the columns of $S$. We denote each copy by $\{a,b,c,d\}$, to calculate $\langle S_{i^{(a)}},S_{i^{(b)}}\rangle\langle S_{i^{(c)}},S_{i^{(d)}}\rangle$:
\begin{align}
 & =\sum_{n_{1}=1}^{m}S_{n_{1}i^{(a)}}S_{n_{1}i^{(b)}}\cdot\sum_{n_{2}=1}^{m}S_{n_{2}i^{(c)}}S_{n_{2}i^{(d)}}\\
 & =\sum_{n_{1}=1}^{m}\left(\sum_{z^{(a)},k^{(a)},z^{(b)},k^{(b)}=1}^{m}g_{k_{1}k_{2}z_{1}}^{(1,a)}g_{k_{1}^{\prime}k_{2}^{\prime}z_{1}^{\prime}}^{(1,b)}\prod_{p=2}^{N-2}g_{z_{p-1}k_{p+1}z_{p}}^{(p,a)}g_{z_{p-1^{\prime}}k_{p+1}^{\prime}z_{p}^{\prime}}^{(p,b)}\cdot g_{z_{N-2}k_{N}\boldsymbol{n}}^{(N-1,a)}g_{z_{N-2}^{\prime}k_{N}^{\prime}\boldsymbol{n}}^{(N-1,b)}\prod_{l=1}^{N}X_{i_{l}^{(a)}k_{l}^{(a)}}\prod_{l=1}^{N}X_{i_{l}^{(b)}k_{l}^{(b)}}\right)\\
 & \cdot\sum_{n_{2}=1}^{m}\left(\sum_{z^{(c)},k^{(c)},z^{(d)},k^{(d)}=1}^{m}g_{k_{1}k_{2}z_{1}}^{(1,c)}g_{k_{1}^{\prime}k_{2}^{\prime}z_{1}^{\prime}}^{(1,d)}\prod_{p=2}^{N-2}g_{z_{p-1}k_{p+1}z_{p}}^{(p,c)}g_{z_{p-1^{\prime}}k_{p+1}^{\prime}z_{p}^{\prime}}^{(p,d)}\cdot g_{z_{N-2}k_{N}\boldsymbol{n}}^{(N-1,c)}g_{z_{N-2}^{\prime}k_{N}^{\prime}\boldsymbol{n}}^{(N-1,d)}\prod_{l=1}^{N}X_{i_{l}^{(c)}k_{l}^{(c)}}\prod_{l=1}^{N}X_{i_{l}^{(d)}k_{l}^{(d)}}\right).\label{eq:full_S_calcultae}
\end{align}

Rearranging and using the independence of the different tensors, we obtain
\[
\mathbb{E}[\langle S_{i^{(a)}},S_{i^{(b)}}\rangle\langle S_{i^{(c)}},S_{i^{(d)}}\rangle]=\sum_{n_{1},n_{2}=1}^{m}\sum_{Z,K}^{m}\prod_{p=1}^{N-1}E\left[g^{(p,a)}g^{(p,b)}g^{(p,c)}g^{(p,d)}\right]\cdot E\left[\prod_{l=1}^{N}\prod_{q\in\{a,b,c,d\}}X_{i_{l}^{(q)}k_{l}^{(q)}}\right]\label{eq:S_4_copies}.
\]
\newpage
\subsection{Replica Patterns and Chains}
\label{sec:replica-patterns}

The variance computation involves products of four replicas of $S$ (Eq.(\ref{eq:S_4_copies})) , indexed by labels $q \in \{a,b,c,d\}$. At each internal bond and at each leaf, the four index values across replicas can collapse in a limited number of ways. We formalize this as follows.

\begin{definition}
\label{def:patterns}
A replica pattern is an equivalence relation in the set $\{i^{(a)}_l,i^{(b)}_l,i^{(c)}_l,i^{(d)}_l\}$ for every $1\le l\le N$. There are four such relations, which we define as    $\mathcal{P} = \{A, U_1, U_2, 4K\}$, where
\begin{align}
4K &:\; \{i^{(a)}_l=i^{(b)}_l=i^{(c)}_l=i^{(d)}_l\}, \label{eq:patterns} \\
A  &:\; \{i^{(a)}_l=i^{(b)}_l\} \,\sqcup\, \{i^{(c)}_l=i^{(d)}_l\}, \nonumber \\
U_1 &:\; \{i^{(a)}_l=i^{(c)}_l\} \,\sqcup\, \{i^{(b)}_l=i^{(d)}_l\}, \nonumber \\
U_2 &:\; \{i^{(a)}_l=i^{(d)}_l\} \,\sqcup\, \{i^{(b)}_l=i^{(c)}_l\}. \nonumber
\end{align}
We refer to $4K$ as the ``four-of-a-kind'' pattern, to $A$ as the ``aligned'' pattern (matching the main-diagonal pairing of the trace), and to $U_1, U_2$ as the two ``unaligned'' pairings. 
We note that patterns $A,U_1,U_2$  act like a 2 pair pattern, they make the tensors act like 2 independent copies of pairs. So we refer to them as $2P$  patterns.

\label{def:Chain}
Each  $i^{(q)},q\in\{a,b,c,d\} $, is an $N$-tuple itself , i.e.,
\begin{align}
    i^{(q)}&=(i^{(q)}_1,i^{(q)}_2,...,i^{(q)}_N)\in [d]^N.\\
\end{align}
Each sub-index $(i^{(a)}_l,i^{(b)}_l,i^{(c)}_l,i^{(d)}_l)\in[d]^4$ induces a pattern $P^l$  for each $ l\in[N]$.
Finally, we define a chain of patterns  $(P_1,...,P_N)$ induced by the index $(i^{(a)},i^{(b)},i^{(c)},i^{(d)})\in[d]^{4\times N}$.
\end{definition}
\pagebreak
\section{Evaluating M}
We want to evaluate each entry in $M$, to do so we construct the following theorem.
\begin{theorem}[Entry-wise decomposition of $M$]
\label{thm:per-entry-M}
Let $S \in \mathbb{R}^{m \times d^N}$ be the matricized sketching tensor network of Eq.~\eqref{eq:s_def}, with i.i.d.\ core entries $G^{(l)}_{ij},\, X^{(l)}_{ij} \sim Q$, where $Q$ is a (real) distribution satisfying
\[
\mathbb{E}[Q] = \mathbb{E}[Q^3] = 0, \qquad \mathbb{E}[Q^2] = \sigma_2, \qquad \mathbb{E}[Q^4] = \sigma_4.
\]
Let $M$ be the matrix defined in Eq.~\eqref{eq:M_def}, and let  $\mathcal{P}$ be the four replica patterns from Def.~\ref{def:patterns}. 
\begin{equation}
\label{eq:per-entry-M}
M\!\left[(i^a,i^c),\,(i^b,i^d)\right]
\;=\;
V(P_1,\dots,P_N)
\end{equation}
where 
\begin{equation}
\label{eq:chain-value}
V(P_1,\dots,P_N)
\;=\;
\bigl\langle u^{P_N}_{\mathrm{end}},
\, W_{P_{N-1}}\, W_{P_{N-2}} \cdots W_{P_{3}}\,
u^{(P_1,P_2)}_{\mathrm{init}}\rangle
\end{equation}
and  $W_P\in \C^{4\times 4}\text{ and } u^{(P_1,P_2)}_{\mathrm{init}},\, u^{P_N}_{\mathrm{end}}\in\C^{4}$   are the Markov chain components.;
\end{theorem}
\subsection{Sketch Of The Proof}
\subsubsection{Proof logic}
Looking at $S$  as a sum (Eq. (\ref{eq:full_S_calcultae)}), we notice that each core tensor $g^{(p)}$ for $p=3,...,N-2$ share indices with a Kronecker tensor $X^{(p+1)}$  (which is a simple Gaussian tensor affected by the choice of indices) and 2 other core tensors $g^{(p-1)},g^{(p+1)}$.
Hence, the core tensors $g^{(p)}$  act like a linear matrix product with an added interaction with $X^{(p+1)}$.
Therefore, we use the approach of  a \emph{Markov like} model where in order to analyze the sum in Eq. (\ref{eq:full_S_calcultae}) for $p$ tensors, we need information derived from $p-1$  tensors.
Our initial state is derived from $g^{(1)}$, we use our model to calculate the count for $N-2$  tensors $g$ , and the end state is the final tensor $g^{(N-1)}$ .
At each step we count all possible permutations of patterns $\mathcal{P}$  for the next tensor and the weight for that permutation given from the expectation of the tensors in that permutation.

\subsubsection{Proof Road-map.}
The remainder of this section is devoted to proving Theorem~\ref{thm:per-entry-M} and identifying its building blocks. The starting point is the observation that $M$ is defined entry-wise in~\eqref{eq:M_def} as a fourth-order moment of $S$, and that $S$ itself is built from two independent random ingredients: the core gadget tensors $g^{(p)}$ (built in turn from the Gaussian matrices $G^{(p,1)}, G^{(p,2)}$, see Eq.~\eqref{eq:g_def}) and the Kronecker leaf matrices $X^{(l)}$.   
The task is therefore to evaluate, for a given $(i^a, i^b, i^c, i^d)$, a single entry of $M$. We carry this out in three steps:
\begin{enumerate}
    \item Local expectation  (section~\ref{sec:local_expectation}) -  Each sub-index $l\in [N]$ is effected by a fixed set of elements, we evaluate all possible expectations gained.
    \item \emph{Combinatorics of non-vanishing terms} (Section~\ref{sec:combinatorics}). The moment assumption $\mathbb{E}[Q] = \mathbb{E}[Q^3] = 0$ forces every non-vanishing term to be included  in $\mathcal{P}$. This reduces the index sum to a sum over all possible pattern chains $(P_1, \dots, P_N) \in \mathcal{P}^N$.
    \item Markov chain  (Section~\ref{sec:markov}) - We use step 1 to count the amount of possible permutations of patterns, and we use step 2 for each permutation to evaluate its expectation.
\end{enumerate}

\subsection{Local Expectation\label{sec:local_expectation}}

\begin{definition}[Meet of patterns]\label{def:meet}
Let $\mathcal{P} = \{A, U_{1}, U_{2}, 4K\}$ and adjoin a fifth value 
$\mathrm{Other}$. For $K_{1}, K_{2} \in \mathcal{P}$ define
\[
K_{1} \wedge K_{2} \;=\;
\begin{cases}
K_{2} & K_{1} = 4K, \\
K_{1} & K_{2} = 4K, \\
K_{1} & K_{1} = K_{2}, \\
\mathrm{Other} & \text{otherwise.}
\end{cases}
\]
We extend the gadget and leaf expectations of Lemmas~\ref{lem:mu-g} 
and~\ref{lem:X-leaf-full} to $\mathrm{Other}$ by the convention
\[
\mu^{g}_{\mathrm{Other},\,\cdot} \;=\; \mu^{g}_{\cdot,\,\mathrm{Other}} 
\;=\; \mu^{X}_{\cdot,\,\mathrm{Other}} \;=\; 0.
\]

As a direct consequence, for each outer leaf pattern $I\in\mathcal{P}$ 
the set of inner patterns with non-zero leaf expectation is
\[
\mathcal{K}(I) \;:=\; \{B\in\mathcal{P}:\mu^{X}_{I,B}\neq 0\} \;=\;
\begin{cases}
\{A,U_{1},U_{2},4K\} & I=4K,\\
\{A,4K\} & I\in\{A,U_{1},U_{2}\},
\end{cases}\label{eq:K_I_remain}
\]
by inspection of the table in Lemma~\ref{lem:X-leaf-full}. We call 
$\mathcal{K}(I)$ the \emph{admissible inner patterns} for some index pattern $I$.
\end{definition}

$K_{1} \wedge K_{2}$ is the $4$-replica pattern of two indices that are forced to agree iff both component patterns agree: 
a $4K$ pattern imposes no constraint, two equal $2P$ patterns reinforce 
each other, and two distinct $2P$ patterns share no common pair, 
yielding $\mathrm{Other}$ (no two replicas coincide). The vanishing 
convention reflects that in this last case the corresponding tensor 
expectation has each replica appearing an odd number of times, hence 
vanishes by $\mathbb{E}[Q] = \mathbb{E}[Q^{3}] = 0$ as in the proofs 
of the lemmas to come~\ref{lem:mu-g} and~\ref{lem:mu-X}.


\begin{lemma}[Pattern multiplicity]\label{lem:N-count}
For a single index of dimension $\widetilde m$ and a pattern $s \in \mathcal{P}$, 
define $N_{s}(\widetilde m)$ to be the number of $4$-tuples in 
$[\widetilde m]^{4}$ realizing $s$. Then
\[
N_{s}(\widetilde m) \;=\;
\begin{cases}
\widetilde m & s = 4K, \\
\widetilde m (\widetilde m - 1) & s \in \{A, U_{1}, U_{2}\}.
\end{cases}
\]
\end{lemma}

\begin{proof}
Given some index $i$  we look at the 4 copies of that index, $i^{(a)},i^{(b)},i^{(c)},i^{(d)}$.
If the 4 indices coincide in a $4K$  pattern, i.e. $$i^{(a)}=i^{(b)}=i^{(c)}=i^{(d)}$$ then we simply fix $i^{(a)}$  that has $\tilde{m}$ options, and the $4K$  pattern enforces all other indices to be the same. All in all, only $\tilde{m}$  options.
If the $4$  indices coincide in one of the $2P$  patterns, w.l.o.g lets assume the pattern is $A$ we get that $$i^{(a)}=i^{(b)}\not=i^{(c)}=i^{(d)}$$
then we fix $i^{(a)}$  which has $\tilde{m}$ possibilities, once chosen this defines $i^{(b)}$. Then, $i^{(c)}$  has $\tilde{m}-1$  possibilities because it must be different then the one chosen for $i^{(a)}$, once chosen it fixes $i^{(d)}$. All in all, $\tilde{m}(\tilde{m}-1)$  possibilities.

\end{proof}

We look again at Eq. (\ref{eq:S_4_copies}):
\[
\sum_{n_{1},n_{2}=1}^{m}\sum_{Z,K}^{m}\prod_{p=1}^{N-1}E\left[g^{(p,a)}g^{(p,b)}g^{(p,c)}g^{(p,d)}\right]\cdot\prod_{l=1}^{N}E\left[\prod_{q\in\{a,b,c,d\}}X_{i_{l}^{(q)}k_{l}^{(q)}}\right].
\]
It consists of two parts:
\begin{itemize}
\item Per index part : $\prod_{l=1}^{N}E\left[\prod_{q\in\{a,b,c,d\}}X_{i_{l}^{(q)}k_{l}^{(q)}}\right]$ - varies for the indices $i^{(a)},i^{(b)},i^{(c)},i^{(d)}$.
\item The common part : $\sum_{n_{1},n_{2}=1}^{m}\sum_{Z,K}^{m}\prod_{p=1}^{N-1}E\left[g^{(p,a)}g^{(p,b)}g^{(p,c)}g^{(p,d)}\right]$ - is fixed for every index and can be written equivalently \begin{equation}
    \sum_{n_{1},n_{2}=1}^{m}\sum_{Z,K=1}^{m}\prod_{q\in\{a,b,c,d\}}\left(g_{k_{1}^{(q)}k_{2}^{(q)}z_{1}^{(q)}}^{(1,q)}\prod_{p=2}^{N-2}g_{z_{p-1}k_{p+1}z_{p}}^{(p,q)}\cdot g_{z_{N-2}^{(q)}k_{N}^{(q)}\boldsymbol{n^{(q)}}}^{(N-1,q)}\right)\label{eq:common}.
\end{equation}
\end{itemize}
We begin by inspecting the expectation of a single value in the equation above:
\begin{align}
\mathbb{E}[g^{(p)}_{k_1k_2z_1}] & =E\left[\sum_{l=1}^{m}G^{(1,p)}_{k_{1}k_{2}^{(1)}l}\cdot G^{(2,p)}_{lk_{i}^{(2)}z_{1}}\right]\\
 & =\sum_{l=1}^{m}\mathbb{E}[G^{(1,p)}_{k_{1}k_{2}^{(1)}l}]\cdot \mathbb{E}[G^{(2,p)}_{lk_{i}^{(2)}z_{1}}]\\
 & =0
\end{align}

Where we used expectation linearity and the assumption of the distribution.
Each summand in Eq. (\ref{eq:common}) consists of a product of 4 copies of $g^{(p)}$. Since $\mathbb{E}[g^{(p)}]=0$, if the relevant indices don't coincide, the entire product vanishes. Therefore, each tensor replicas $g^{(p)}$ with $p\in [N-1]$  must correspond to a pattern in $\mathcal{P}$.  We evaluate the expectation received from the $g^{(p)}$  tensors in the following lemma.
\begin{lemma}\label{lem:mu-g}
Let
\[
g^{(p)}_{z_{p-1}\,k^{(1)}_{p+1}\,k^{(2)}_{p+1}\,z_{p}}
\;=\;
\sum_{l=1}^{m}\,
G^{(p,1)}_{z_{p-1}\,k^{(1)}_{p+1}\,l}\;
G^{(p,2)}_{l\,k^{(2)}_{p+1}\,z_{p}}
\]
where each entry in  $G^{(p,1)},G^{(p,2)}$  is drawn i.i.d.\ from
a real distribution $Q$ satisfying
\[
\mathbb{E}[Q]=\mathbb{E}[Q^{3}]=0,\qquad
\mathbb{E}[Q^{2}]=\sigma_{2},\qquad
\mathbb{E}[Q^{4}]=\sigma_{4}.
\]
We fix the $4$-replica patterns of the gadget's two halves:
\[
L\in\{A,U_{1},U_{2},4K\}\ \text{ for the left-half outer indices }(z_{p-1},\,k^{(1)}_{p+1}) \text{ given from $G^{(p,1)}$},
\]
\[
R\in\{A,U_{1},U_{2},4K\}\ \text{ for the right-half outer indices }(k^{(2)}_{p+1},\,z_{p})\text{ given from $G^{(p,2)}$}.
\]
Then
\[
\mu^{g}_{L,R}\ :=\ \mathbb{E}\!\left[\prod_{q\in\{a,b,c,d\}}g^{(p,q)}\right]
\]
is given by the $4\times 4$ table
\[
\renewcommand{\arraystretch}{1.4}
\begin{array}{c|cccc}
L\,\backslash\,R & A & U_{1} & U_{2} & 4K \\\hline
A     & m^{2}\sigma_{2}^{4}
      & m\sigma_{2}^{4}
      & m\sigma_{2}^{4}
      & m(m-1)\sigma_{2}^{4}+m\sigma_{4}\sigma_{2}^{2} \\
U_{1} & m\sigma_{2}^{4}
      & m^{2}\sigma_{2}^{4}
      & m\sigma_{2}^{4}
      & m(m-1)\sigma_{2}^{4}+m\sigma_{4}\sigma_{2}^{2} \\
U_{2} & m\sigma_{2}^{4}
      & m\sigma_{2}^{4}
      & m^{2}\sigma_{2}^{4}
      & m(m-1)\sigma_{2}^{4}+m\sigma_{4}\sigma_{2}^{2} \\
4K    & m(m-1)\sigma_{2}^{4}+m\sigma_{4}\sigma_{2}^{2}
      & m(m-1)\sigma_{2}^{4}+m\sigma_{4}\sigma_{2}^{2}
      & m(m-1)\sigma_{2}^{4}+m\sigma_{4}\sigma_{2}^{2}
      & 3m(m-1)\sigma_{2}^{4}+m\sigma_{4}^{2} \\
\end{array}
\]
\end{lemma}
\begin{proof}
The two halves $G^{(p,1)},G^{(p,2)}$ are independent and share only the
"dummy" index~$l$. Substituting the definition of $g^{(p,q)}$ eq.(\eqref{eq:g_def}) and using the independence, we obtain
\begin{equation}\label{eq:mug-raw}
\mu^{g}_{L,R}\;=\!\!\sum_{l_{a},l_{b},l_{c},l_{d}=1}^{m}\!\!
\mathbb{E}\!\Big[\prod_{q}G^{(p,1),q}_{z_{p-1}^{(q)}\,k^{(1),q}_{p+1}\,l_{q}}\Big]\;
\mathbb{E}\!\Big[\prod_{q}G^{(p,2),q}_{l_{q}\,k^{(2),q}_{p+1}\,z_{p}^{(q)}}\Big].
\end{equation}

There are $m^{4}$ tuples of $(l_{a},l_{b},l_{c},l_{d})$. As stated above, only tuples that vanish are of pattern $4K$  or $2P$. 
We denote $s\in\{4K, A, U_1,U_2\}$ as the pattern type for $(l_{a},l_{b},l_{c},l_{d})$. All other tuples have at least one $l_{q}$  that has an odd number of repetitions. We denote $N_s$  like in lemma (\ref{lem:N-count}).
Splitting the sum from Eq.  (~\eqref{eq:mug-raw}) along these patterns, we obtain
\begin{equation}\label{eq:mug-split}
\mu^{g}_{L,R}\;=\!\!\sum_{s\in\{A,U_{1},U_{2},4K\}}\!\!
N_{s}\cdot
\underbrace{\mathbb{E}\!\Big[\prod_{q}G^{(p,1),q}_{z_{p-1}^{(q)}\,k^{(1),q}_{p+1}\,l_{q}}\Big]}_{\text{value }w_{L,s}}
\cdot
\underbrace{\mathbb{E}\!\Big[\prod_{q}G^{(p,2),q}_{l_{q}\,k^{(2),q}_{p+1}\,z_{p}^{(q)}}\Big]}_{\text{value }w_{R,s}}
\;+\;0,
\end{equation}
where  $0$ accounts for tuples in a pattern that vanish, and $w_{L,s}, w_{R,s}$ are the per-half expectations, which depend only on the joint pattern realised by $(z_{p-1},k^{(1)}_{p+1},l)$ and
$(l,k^{(2)}_{p+1},z_{p})$ respectively.

\paragraph{Why tuples outside the four patterns vanish.}
For each tensor $G^{(p,1),q}$, its three indices
$(z_{p-1}^{(q)},k^{(1),q}_{p+1},l_{q})$ define a triple. The four
replicas yield four such triples. We group replicas by equal triples; if
any group has odd size, its contribution to the expectation is a product
involving an odd power of a centred variable, which vanishes by
$\mathbb{E}[G]=\mathbb{E}[G^{3}]=0$. Hence, non-vanishing requires the
joint replica partition to be in $\{A,U_{1},U_{2},4K\}$. The same
applies to $G^{(p,2)}$. In particular, the joint $l$-partition must lie
in this set as well.

\paragraph{Computing $w_{L,s}$.}
For a fixed $L$ and $s$, the triple $(z_{p-1}, k^{(1)}_{p+1}, l)$ joins 
the outer pattern $L$ with the $l$-pattern $s$ component-wise. The 
joint pattern is the meet $L \wedge s$ (cf.\ Definition~\ref{def:meet}): 
two replicas $q, q'$ index the same entry of $G^{(p,1)}$ iff they 
coincide on the outer indices \emph{and} on $l$, i.e.\ iff they are in 
the same block of $L \wedge s$.
\begin{itemize}
  \item If $L \wedge s = \mathrm{Other}$, the four replicas index four 
        distinct entries (no coincidences). The expectation factors as 
        a product of four single-$G$ moments, each equal to 
        $\mathbb{E}[Q] = 0$, hence $w_{L,s} = 0$.
  \item If $L \wedge s \in \{A, U_{1}, U_{2}\}$, the four replicas split 
        into two pairs sharing entries, with the two pairs independent. 
        The expectation factors as $\mathbb{E}[Q^{2}] \cdot \mathbb{E}[Q^{2}] = \sigma_{2} \cdot \sigma_{2}$, hence $w_{L,s} = \sigma_{2}^{2}$.
  \item If $L \wedge s = 4K$, all four replicas index the same entry, 
        contributing $\mathbb{E}[Q^{4}] = \sigma_{4}$, hence $w_{L,s} = \sigma_{4}$.
\end{itemize}
The same rules give $w_{R,s}$ from $R \wedge s$.

\paragraph{When is $w_{L,s}$ non-zero?}
By the previous paragraph, $w_{L,s} \ne 0$ iff $L \wedge s \ne \mathrm{Other}$. 
Inspecting Definition~\ref{def:meet}, the set of $s \in \mathcal{P}$ satisfying 
this is
\[
\mathcal{L}_{L} \;:=\; \{s \in \mathcal{P} : L \wedge s \ne \mathrm{Other}\}
\;=\;
\begin{cases}
\{L, 4K\} & L \in \{A, U_{1}, U_{2}\}, \\
\mathcal{P} & L = 4K,
\end{cases}
\]
and similarly for $\mathcal{L}_{R}$. The only non-vanishing terms in 
the sum~\eqref{eq:mug-split} are those with $s \in \mathcal{L}_{L} \cap \mathcal{L}_{R}$.

\paragraph{Reading off the four cases.}
For each $(L,R)$ we list the "surviving" $s\in\mathcal L_{L}\cap\mathcal L_{R}$ and the half values:

\smallskip\noindent\emph{Case 1: $L=R\in\{A,U_{1},U_{2}\}$.}
$\mathcal L_{L}\cap\mathcal L_{R}=\{L,4K\}$ and $w_{L,s}=w_{R,s}=\sigma_{2}^{2}$ for all $s$:
\[
\mu^{g}_{L,L}=N_{L}\sigma_{2}^{4}+N_{4K}\sigma_{2}^{4}
=\bigl[m(m-1)+m\bigr]\sigma_{2}^{4}=m^{2}\sigma_{2}^{4}.
\]

\smallskip\noindent\emph{Case 2: $L\ne R$, both in $\{A,U_{1},U_{2}\}$.}
$\mathcal L_{L}\cap\mathcal L_{R}=\{4K\}$ only. We obtain $w_{L,s}=w_{R,s}=\sigma_{2}^{2}$
\[
\mu^{g}_{L,R}=N_{4K}\sigma_{2}^{4}=m\sigma_{2}^{4}.
\]

\smallskip\noindent\emph{Case 3: exactly one of $L,R$ is $4K$.}
W.L.G $L=4K,\ R\in\{A,U_{1},U_{2}\}$ and we get$\mathcal L_{L}\cap\mathcal L_{R}=\{R,4K\}$.
\begin{itemize}\setlength\itemsep{1pt}
  \item $s=R$: meet on left and right are a $2P$, so $w_{L,R}=w_{R,R}=\sigma_{2}^{2}$. Contribution: $N_{R}\sigma_{2}^{4}=m(m-1)\sigma_{2}^{4}$.
  \item $s=4K$: meet on left is $4K\wedge 4K=4K$, so $w_{L,4K}=\sigma_{4}$. Meet on right is $R\wedge 4K=R$, so $w_{R,4K}=\sigma_{2}^{2}$. Contribution: $N_{4K}\sigma_{4}\sigma_{2}^{2}=m\sigma_{4}\sigma_{2}^{2}$.
\end{itemize}
\[
\mu^{g}_{4K,R}=\mu_{L,4K}^g=m(m-1)\sigma_{2}^{4}+m\sigma_{4}\sigma_{2}^{2}.
\]

\smallskip\noindent\emph{Case 4: $L=R=4K$.}
$\mathcal L_{L}\cap\mathcal L_{R}=\{A,U_{1},U_{2},4K\}$.
\begin{itemize}\setlength\itemsep{1pt}
  \item $s\in\{A,U_{1},U_{2}\}$: meet on each side is $4K\wedge s=s$ (a $2P$), so both halves give $\sigma_{2}^{2}$. Total contribution: $3\cdot N_{A}\sigma_{2}^{4}=3m(m-1)\sigma_{2}^{4}$.
  \item $s=4K$: meet on each side is $4K$, so both halves give $\sigma_{4}$. Contribution: $N_{4K}\sigma_{4}^{2}=m\sigma_{4}^{2}$.
\end{itemize}
\[
\mu^{g}_{4K,4K}=3m(m-1)\sigma_{2}^{4}+m\sigma_{4}^{2}.
\]

The four cases reproduce the table.
\end{proof}
Next we calculate the expectation of the index part from equation Eq. (\ref{eq:S_4_copies})
And hence the following lemma:\begin{lemma}\label{lem:X-leaf-full}
Let $X^{(p)} \in \mathbb{R}^{m \times d}$ be the leaf tensor with entries 
$X^{(p)}_{i^{(p)}k^{(p)}}$ i.i.d.\ from a real distribution $Q$ satisfying 
$\mathbb{E}[Q] = \mathbb{E}[Q^3] = 0,\ \mathbb{E}[Q^2] = \sigma_2,\ 
\mathbb{E}[Q^4] = \sigma_4$. Fix the $4$-replica patterns of the two indices:
\[
I \in \{A, U_1, U_2, 4K\} \text{ for the outer (physical) index } i_l, \qquad
B \in \{A, U_1, U_2, 4K\} \text{ for the inner (bond) index } k_l.
\]
Then
\[
\mu^X_{I,B} \;:=\; \mathbb{E}\!\left[\,\prod_{q\in\{a,b,c,d\}} 
X^{(l),q}_{\,i^{(q)}_l,\,k^{(q)}_l}\right]
\]
is given by the $4\times 4$ table
\[
\begin{array}{c|cccc}
I\,\backslash\,B & A & U_1 & U_2 & 4K \\ \hline
A   & \sigma_2^2 & 0          & 0          & \sigma_2^2 \\
U_1 & 0          & \sigma_2^2 & 0          & \sigma_2^2 \\
U_2 & 0          & 0          & \sigma_2^2 & \sigma_2^2 \\
4K  & \sigma_2^2 & \sigma_2^2 & \sigma_2^2 & \sigma_4
\end{array}
\]
Equivalently, with $\wedge$ denoting the meet (componentwise refinement) of equivalence relations,
\[
\mu^X_{I,B} \;=\;
\begin{cases}
\sigma_4    & I\wedge B = 4K,\\
\sigma_2^2  & I\wedge B \in \{A,U_1,U_2\},\\
0           & \text{otherwise (meet contains an odd-sized block).}
\end{cases}
\]
\end{lemma}

\begin{proof}
Two replicas $q,q'\in\{a,b,c,d\}$ index the same entry of $X^{(l)}$ iff 
they agree on both $i_l$ and $k_l$, i.e.\ iff $q\sim_I q'$ and $q\sim_B q'$.
We compute $\mu^X_{I,B}$ by enumerating the cases.

\smallskip\noindent\textbf{Case $I=B=4K$.} All four replicas share the same 
entry, so $\mu^X_{4K,4K} = \mathbb{E}[Q^4] = \sigma_4$.

\smallskip\noindent\textbf{Case $I=4K,\ B\in\{A,U_1,U_2\}$.} The outer index 
imposes no constraint, so the replica coincidences are exactly those of $B$: 
the four replicas split into two pairs according to $B$. Each pair 
contributes $\mathbb{E}[Q^2]=\sigma_2$, hence $\mu^X_{4K,B}=\sigma_2^2$. The 
symmetric case $I\in\{A,U_1,U_2\},\ B=4K$ is identical.

\smallskip\noindent\textbf{Case $I=B\in\{A,U_1,U_2\}$.} Both indices induce 
the same pairing of replicas, so the four replicas split into two pairs 
sharing the same entry. Each pair contributes $\sigma_2$, hence 
$\mu^X_{I,I}=\sigma_2^2$.

\smallskip\noindent\textbf{Case $I,B\in\{A,U_1,U_2\}$ with $I\ne B$.} 
Any two distinct $2P$ partitions ($A,U_1,U_2$) share no common pair. Hence no two replicas index the same entry, and the four factors are independent:
\[
\mu^X_{I,B} \;=\; \prod_{q\in\{a,b,c,d\}} \mathbb{E}\!\left[X^{(l),q}_{i^{(q)}_l,k^{(q)}_l}\right]
\;=\; \mathbb{E}[Q]^4 \;=\; 0.
\]

\smallskip
Collecting the four cases yields the table.
\end{proof}

\subsection{Combinatorics}\label{sec:combinatorics}

Lemmas~\ref{lem:mu-g} and~\ref{lem:mu-X} give the per-site per-pattern 
expectations of the gadget and leaf tensors. To assemble these into a 
transfer matrix we count, at each site, the number of index 
configurations realizing each pattern and weight them by their 
expectations. Each site $p\in\{3,\dots,N-2\}$ only needs to know how 
site $p-1$ ended in order to determine how the shared bond acts: given 
the bond pattern emitted by site $p-1$, we enumerate the legal patterns 
at site $p$ and weight each by its expectation. This is the Markov 
structure formalized in Theorem~\ref{thm:W-bulk} below.
\subsubsection{Markov-like chains}\label{sec:markov}

We now assemble the per-site weights into a transfer matrix. The 
construction is \emph{Markov-like}: because adjacent gadgets share only 
the bond index $z_{p-1}$, the joint contribution of sites 
$1, \dots, p$ depends on sites $1, \dots, p-1$ only through the pattern 
of that shared bond.
\subsection{Transfer Matrix Formula}
\begin{theorem}[Transfer matrix formula]\label{thm:W-bulk}
Fix a gadget site $p \in \{3, \dots, N-2\}$ and a pattern $P_{p+1} \in \mathcal{P}$ for the outer index of the leaf  $X^{(p+1)}$ attached to $g^{(p)}$ through the shared index $k_{p+1}$ 
(cf.~\eqref{eq:g_def}). 
There exists a $4 \times 4$ matrix $W^{P_{p+1}}$, indexed by bond patterns $Z_{\mathrm{out}}, Z_{\mathrm{in}} \in \mathcal{P}$ in the basis 
$(A, U_{1}, U_{2}, 4K)$, such that the joint contribution of gadget $g^{(p)}$ together with leaf $X^{(p+1)}$ to the sum~\eqref{eq:S_4_copies}, conditional on the incoming bond pattern $Z_{\mathrm{in}}$ of $z_{p-1}$ and the outgoing bond pattern $Z_{\mathrm{out}}$ of $z_{p}$, is the entry 
$W^{P_{p+1}}_{Z_{\mathrm{out}}, Z_{\mathrm{in}}}$. Moreover,
\begin{equation}\label{eq:W-bulk}
W^{P_{p+1}}_{Z_{\mathrm{out}}, Z_{\mathrm{in}}}
\;=\;
N_{Z_{\mathrm{out}}}(m)
\!\!\sum_{K \,\in\, \mathcal{K}(P_{p+1})}\!\! \mu^{X}_{P_{p+1}, K}
\!\!\!\!\sum_{\substack{K_{1}, K_{2} \in \mathcal{P} \\ K_{1} \wedge K_{2} \,=\, K}}\!\!\!\!
N_{K_{1}}(m_{1})\, N_{K_{2}}(m_{2})\,
\mu^{g}_{\,Z_{\mathrm{in}} \wedge K_{1},\, K_{2} \wedge Z_{\mathrm{out}}}.
\end{equation}
\end{theorem}
\begin{proof}
\textbf{Setting up the model.} 
We want to count and weigh all possible permutation for a chain sized $p$. We look at the site $p$ of our network, which consists of $g^{(p)} \text{ and } X^{(p+1)}$, and note that we have 5 new indices from the to sum on, each with their own restriction. 
We remember the indices for $g_{z_{p-1}k_{p+1}^{(1)}k_{p+1}^{(2)}z_{p}}^{(p)}$, eq.(\ref{eq:g_def}) and $X^{(p+1)}_{k_{p+1},i_{p+1}}$, eq.(\ref{eq:S_4_copies}). 
Meaning we have 5 new indices with the following restrictions:
\begin{enumerate}
    \item $z_{p-1}$  is shared with $g^{(p-1)}$, so we assume it is already fixed with a pattern $Z_{in}\in \mathcal{P}$   by the previous $p-1$  sized chain.
    \item $k_{p+1}^{(1)}k_{p+1}^{(2)}$  are shared between $g^{(p)}$  and $X^{(p+1)}$. Because $X^{(p+1)}$  is directly connected to the choice of index $i_{p+1}$ makes the $k$  indices act in a cretin way.
    \item $z_p$  is shared with the next tensor, $g^{(p+1)}$, so we remember its pattern $Z_{out}\in\mathcal{P}$ .
\end{enumerate}

With all of this in mind, we construct a matrix $W^{(P_{p+1})}\in \R^{4\times 4}$ for our site, on the basis $\mathcal{P}=\{4K,A,U_1,U_2\}$. Where each element is the weight and count of the non vanishing permutation starting with a pattern $Z_{in}$  and ending with $Z_{out}$. The value itself is based on the lemmas from earlier \ref{lem:mu-g} and~\ref{lem:X-leaf-full}.

\smallskip\noindent\textbf{Gadget layer.} 
The gadget $g^{(p)}$ is built from two halves, see Eq.~\ref{eq:g_def}, 
$G^{(p,1)}$ and $G^{(p,2)}$. The left half $G^{(p,1)}$ carries the 
indices $(z_{p-1}, k^{(1)}_{p+1})$ and the right half $G^{(p,2)}$ 
carries $(k^{(2)}_{p+1}, z_{p})$. Two replicas agree on a half iff 
they agree on both of its outer indices, so the possible patterns of the two 
halves are
\[
L \;=\; Z_{\mathrm{in}} \wedge K_{1}, \qquad 
R \;=\; K_{2} \wedge Z_{\mathrm{out}},
\]
where $K_{1}, K_{2} \in \mathcal{P}$ are the patterns of the two factors of $k_{p+1}$. 
For fixed $(K_{1}, K_{2})$ by Lemma~\ref{lem:N-count} there are $N_{K_{1}}(m_{1})\, N_{K_{2}}(m_{2})$ value-tuples for the factors, and by Lemma~\ref{lem:mu-g} they contribute the gadget weight 
$\mu^{g}_{L, R}$. Clashes are absorbed automatically: if $L = \mathrm{Other}$ or $R = \mathrm{Other}$, the gadget factor vanishes by the convention of Definition~\ref{def:meet}.
We get that for some fixed $(K_1,K_2)$  the gadget $g^{(p)}$ adds $$N_{K_{1}}(m_{1})\, N_{K_{2}}(m_{2})\,
\mu^{g}_{\,Z_{\mathrm{in}} \wedge K_{1},\, K_{2} \wedge Z_{\mathrm{out}}}$$
\smallskip\noindent\textbf{Leaf layer (middle).} 
Given that the index $i_{(p+1)}$  of tensor $X^{(p+1)}$ is of a pattern $P_{p+1}$ we get that the values that arise from the replicas of $X^{(p+1)}$ are non-vanishing only when $K\in\mathcal{K}(P_{p+1})$, this can be seen by remembering the def (\ref{def:meet})  with the lemma (\ref{lem:X-leaf-full}).
Now for a given $K=K_1\land K_2$  we get that the leaf $X^{(p+1)}$  adds a weight of $\mu_{P_{p+1},K}^X$   as can be seen in lemma (\ref{lem:X-leaf-full})

With both of these paragraphs in mind, we get that the restrictions on $z_{p-1},i_{p+1},k_{p+1}$  induce the formula $$\sum_{K \,\in\, \mathcal{K}(P_{p+1})}\!\! \mu^{X}_{P_{p+1}, K}
\!\!\!\!\sum_{\substack{K_{1}, K_{2} \in \mathcal{P} \\ K_{1} \wedge K_{2} \,=\, K}}\!\!\!\!
N_{K_{1}}(m_{1})\, N_{K_{2}}(m_{2})\,
\mu^{g}_{\,Z_{\mathrm{in}} \wedge K_{1},\, K_{2} \wedge Z_{\mathrm{out}}}$$

\smallskip\noindent\textbf{Outgoing-bond layer (outermost).} 
Finally, we do this for every possible independent index $z_{p}$  which is only restricted by the pattern we are trying to calculate for, $Z_{out}$  so there are $N_{Z_{out}}(m)$  options. Meaning we get the final element formula (\ref{eq:W-bulk})
\end{proof}

\subsubsection{Explicit Transition Matrices}
From the theorem we can write explicitly all 4 cases.
\noindent In the basis $(A, U_{1}, U_{2}, 4K)$ with $m_{1} m_{2} = m$, define the auxiliary 
polynomials
\begin{align*}
\alpha       &\;:=\; (m-1)\,\sigma_{2}^{2} + \sigma_{4}, 
\\ \beta        &\;:=\; (m_{1}+m_{2}-2)\,\sigma_{2}^{2} + \sigma_{4}, \\
\beta_{m_{1}} &\;:=\; (m+m_{1}-2)\,\sigma_{2}^{2} + \sigma_{4},
\\ \beta_{m_{2}} &\;:=\; (m+m_{2}-2)\,\sigma_{2}^{2} + \sigma_{4}, \\
\rho_{m_{1}}  &\;:=\; (m\,m_{1}-1)\,\sigma_{2}^{2} + \sigma_{4},
\\ \rho_{m_{2}}  &\;:=\; (m\,m_{2}-1)\,\sigma_{2}^{2} + \sigma_{4}, \\[4pt]
\Gamma(a,b)  &\;:=\; \bigl(m(m-1) + 2a - b - 1\bigr)\sigma_{2}^{4}
              + (m+b-2)\,\sigma_{2}^{2}\sigma_{4} + \sigma_{4}^{2}, \\
\Theta       &\;:=\; \bigl(m(m+2) - (m_{1}+m_{2}) - 1\bigr)\sigma_{2}^{4}
              + (m_{1}+m_{2}-2)\,\sigma_{2}^{2}\sigma_{4} + \sigma_{4}^{2}, \\
\Delta       &\;:=\; \bigl(3 m(m-1) - (m_{1}+m_{2}-2)\bigr)\sigma_{2}^{6}
              + 3(m+m_{1}+m_{2}-3)\,\sigma_{2}^{4}\sigma_{4} + \sigma_{4}^{3}.
\end{align*}
\[
W^{4K} \;=\; m^{3}\,\sigma_{2}^{2}\,
\begin{pmatrix}
m(m-1)\,\sigma_{2}^{2}\,\alpha & (m-1)\,\sigma_{2}^{2}\,\beta   & (m-1)\,\sigma_{2}^{2}\,\beta   & (m-1)\,\Gamma(m_{1},m_{2}) \\[2pt]
(m-1)\,\sigma_{2}^{2}\,\beta   & m(m-1)\,\sigma_{2}^{2}\,\alpha & (m-1)\,\sigma_{2}^{2}\,\beta   & (m-1)\,\Gamma(m_{1},m_{2}) \\[2pt]
(m-1)\,\sigma_{2}^{2}\,\beta   & (m-1)\,\sigma_{2}^{2}\,\beta   & m(m-1)\,\sigma_{2}^{2}\,\alpha & (m-1)\,\Gamma(m_{1},m_{2}) \\[2pt]
\Gamma(m_{2},m_{1})            & \Gamma(m_{2},m_{1})            & \Gamma(m_{2},m_{1})            & \Delta
\end{pmatrix}.
\]

\[
W^{A} \;=\; m^{3}\,\sigma_{2}^{2}\,
\begin{pmatrix}
m^{2}(m-1)\,\sigma_{2}^{4} & m_{2}(m-1)\,\sigma_{2}^{4} & m_{2}(m-1)\,\sigma_{2}^{4} & m_{2}(m-1)\,\sigma_{2}^{2}\,\rho_{m_{1}} \\[2pt]
m_{1}(m-1)\,\sigma_{2}^{4} & m(m-1)\,\sigma_{2}^{4}    & (m-1)\,\sigma_{2}^{4}      & (m-1)\,\sigma_{2}^{2}\,\beta_{m_{1}}     \\[2pt]
m_{1}(m-1)\,\sigma_{2}^{4} & (m-1)\,\sigma_{2}^{4}     & m(m-1)\,\sigma_{2}^{4}    & (m-1)\,\sigma_{2}^{2}\,\beta_{m_{1}}     \\[2pt]
m_{1}\,\sigma_{2}^{2}\,\rho_{m_{2}} & \sigma_{2}^{2}\,\beta_{m_{2}}         & \sigma_{2}^{2}\,\beta_{m_{2}}         & \Theta
\end{pmatrix}.
\]
\[
W^{U_{1}} \;=\; m^{3}\,\sigma_{2}^{2}\,
\begin{pmatrix}
m(m-1)\,\sigma_{2}^{4}      & m_{1}(m-1)\,\sigma_{2}^{4} & (m-1)\,\sigma_{2}^{4}      & (m-1)\,\sigma_{2}^{2}\,\beta_{m_{1}}     \\[2pt]
m_{2}(m-1)\,\sigma_{2}^{4} & m^{2}(m-1)\,\sigma_{2}^{4} & m_{2}(m-1)\,\sigma_{2}^{4} & m_{2}(m-1)\,\sigma_{2}^{2}\,\rho_{m_{2}} \\[2pt]
(m-1)\,\sigma_{2}^{4}      & m_{1}(m-1)\,\sigma_{2}^{4} & m(m-1)\,\sigma_{2}^{4}     & (m-1)\,\sigma_{2}^{2}\,\beta_{m_{1}}     \\[2pt]
\sigma_{2}^{2}\,\beta_{m_{2}}         & m_{1}\,\sigma_{2}^{2}\,\rho_{m_{2}} & \sigma_{2}^{2}\,\beta_{m_{2}}         & \Theta
\end{pmatrix}.
\]

\[
W^{U_{2}} \;=\; m^{3}\,\sigma_{2}^{2}\,
\begin{pmatrix}
m(m-1)\,\sigma_{2}^{4}      & (m-1)\,\sigma_{2}^{4}      & m_{1}(m-1)\,\sigma_{2}^{4} & (m-1)\,\sigma_{2}^{2}\,\beta_{m_{1}}     \\[2pt]
(m-1)\,\sigma_{2}^{4}      & m(m-1)\,\sigma_{2}^{4}     & m_{1}(m-1)\,\sigma_{2}^{4} & (m-1)\,\sigma_{2}^{2}\,\beta_{m_{1}}     \\[2pt]
m_{2}(m-1)\,\sigma_{2}^{4} & m_{2}(m-1)\,\sigma_{2}^{4} & m^{2}(m-1)\,\sigma_{2}^{4} & m_{2}(m-1)\,\sigma_{2}^{2}\,\rho_{m_{2}} \\[2pt]
\sigma_{2}^{2}\,\beta_{m_{2}}         & \sigma_{2}^{2}\,\beta_{m_{2}}         & m_{1}\,\sigma_{2}^{2}\,\rho_{m_{2}} & \Theta
\end{pmatrix}.
\]


\begin{theorem}[Initial vector]\label{thm:u-init}
With $N \ge 3$ and patterns $P_{1}, P_{2} \in \mathcal{P}$ be the first 2 for the indices of the leaves $X^{(1)}$ and $X^{(2)}$.
For each of the $16$ choices of $(P_{1}, P_{2}) \in \mathcal{P}^{2}$, the joint 
contribution of $X^{(1)}, X^{(2)}$, and the first gadget $g^{(1)}$ to 
the sum~\eqref{eq:S_4_copies} is a vector 
$u^{(P_{1}, P_{2})}_{\mathrm{init}} \in \mathbb{R}^{4}$ indexed by the outgoing bond pattern $Z_{\mathrm{out}} \in \mathcal{P}$ of $z_{1}$, 
whose entries are explicit polynomials in $(m, m_{1}, m_{2}, \sigma_{2}, 
\sigma_{4})$ as by the formula:
\begin{equation}\label{eq:u-init}
\bigl[u^{(P_{1}, P_{2})}_{\mathrm{init}}\bigr]_{Z_{\mathrm{out}}}
\;=\;
N_{Z_{\mathrm{out}}}(m)
\!\!\!\!\!\sum_{\substack{K \,\in\, \mathcal{K}(P_{1}) \\ J \,\in\, \mathcal{K}(P_{2})}}\!\!\!\!\!
N_{K}(m)\, \mu^{X}_{P_{1}, K}\, \mu^{X}_{P_{2}, J}
\!\!\!\!\!\sum_{\substack{K_{1}, K_{2} \,\in\, \mathcal{P} \\ K_{1} \wedge K_{2} \,=\, J}}\!\!\!\!\!
N_{K_{1}}(m_{1})\, N_{K_{2}}(m_{2})\,
\mu^{g}_{\,K \wedge K_{1},\, K_{2} \wedge Z_{\mathrm{out}}}.
\end{equation}


\end{theorem}
\begin{proof}
\textbf{Setting up the model.} 
We want to count and weigh all possible permutations for the start of our chain, i.e. for $p=1,2$. We look at sites 1 and 2 of our network, which consist of $X^{(1)}, g^{(1)}$ and $X^{(2)}$, and this time we have 6 indices to sum on.
We remember the indices for $g_{k_{1}k_{2}^{(1)}k_{2}^{(2)}z_{1}}^{(1)}$, eq.(\ref{eq:g_def}), and the leaves $X^{(1)}_{k_{1},i_{1}}, X^{(2)}_{k_{2},i_{2}}$, eq.(\ref{eq:S_4_copies}).
Note that, unlike interior gadgets, the left outer index of $g^{(1)}$ is the full unfactored inner index $k_{1}$ of $X^{(1)}$, of dimension $m$.
We have the following restrictions:
\begin{enumerate}
    \item $i_{1}, i_{2}$ are the outer indices of $X^{(1)}, X^{(2)}$, so we assume they are already fixed with patterns $P_{1}, P_{2} \in \mathcal{P}$.
    \item $k_{1}$ is shared between $g^{(1)}$ and $X^{(1)}$. Because $X^{(1)}$ is directly connected to the choice of index $i_{1}$ makes the $k_{1}$ index act in a certain way.
    \item $k_{2}^{(1)}k_{2}^{(2)}$ are shared between $g^{(1)}$ and $X^{(2)}$. Because $X^{(2)}$ is directly connected to the choice of index $i_{2}$ makes the $k$ indices act in a certain way.
    \item $z_{1}$ is shared with the next tensor, $g^{(2)}$, so we remember its pattern $Z_{out}\in\mathcal{P}$.
\end{enumerate}

\smallskip\noindent\textbf{Gadget layer.} 
Exactly like the previous theorem. The gadget yields the following counting and weight $$
N_{K_{1}}(m_{1})\, N_{K_{2}}(m_{2})\,
\mu^{g}_{\,K \wedge K_{1},\, K_{2} \wedge Z_{\mathrm{out}}}$$

\smallskip\noindent\textbf{Leaf layer (middle).} 
Given that the index $i_{1}$ of tensor $X^{(1)}$ is of a pattern $P_{1}$, we get that the values that arise from the replicas of $X^{(1)}$ are non-vanishing only when $K\in\mathcal{K}(P_{1})$, this can be seen by remembering the def (\ref{def:meet}) with the lemma (\ref{lem:X-leaf-full}). For such a $K$, the leaf $X^{(1)}$ adds a weight of $\mu^{X}_{P_{1},K}$.
Similarly, given that the index $i_{2}$ of tensor $X^{(2)}$ is of a pattern $P_{2}$, the replicas of $X^{(2)}$ are non-vanishing only when $J := K_{1} \wedge K_{2} \in \mathcal{K}(P_{2})$, and for such $J$ the leaf $X^{(2)}$ adds a weight of $\mu^{X}_{P_{2},J}$.
Moreover, we note that only the $J$  pattern affects both sides of $g^{(1)}$, as seen in the inner sum.
We get that the restrictions on $i_{1},i_{2},k_{1},k_{2}$ induce the formula $$\!\!\!\!\!\sum_{\substack{K \,\in\, \mathcal{K}(P_{1}) \\ J \,\in\, \mathcal{K}(P_{2})}}\!\!\!\!\!
N_{K}(m)\, \mu^{X}_{P_{1}, K}\, \mu^{X}_{P_{2}, J}
\!\!\!\!\!\sum_{\substack{K_{1}, K_{2} \,\in\, \mathcal{P} \\ K_{1} \wedge K_{2} \,=\, J}}\!\!\!\!\!
N_{K_{1}}(m_{1})\, N_{K_{2}}(m_{2})\,
\mu^{g}_{\,K \wedge K_{1},\, K_{2} \wedge Z_{\mathrm{out}}}$$

\smallskip\noindent\textbf{Outgoing-bond layer (outermost).} 
Again, we do this for every possible independent index $z_{1}$ restricted by the pattern $Z_{out}$ with $N_{Z_{out}}(m)$ options. Meaning we get the final element formula (\ref{eq:u-init})
\end{proof}
\subsubsection{Explicit starting state}
\paragraph{Symmetry.} 
Every initial vector satisfies the $U_{1} \leftrightarrow U_{2}$ symmetry
\[
\bigl[u^{(P_{1}, P_{2})}_{\mathrm{init}}\bigr]_{Z}
\;=\;
\bigl[u^{(\sigma P_{1}, \sigma P_{2})}_{\mathrm{init}}\bigr]_{\sigma Z},
\]
where $\sigma : \mathcal{P} \to \mathcal{P}$ swaps $U_{1}$ and $U_{2}$ and fixes $A, 4K$. 
It therefore suffices to list the $10$ orbit representatives below; the remaining $6$ 
vectors follow by applying $\sigma$ to both leaf patterns and the bond pattern. For the 
$4$ self-symmetric representatives ($P_{1}, P_{2} \in \{A, 4K\}$), the symmetry forces 
$[u]_{U_{1}} = [u]_{U_{2}}$.

\paragraph{$u^{(A, A)}_{\mathrm{init}}$.}
\begin{align*}
[u^{(A, A)}_{\mathrm{init}}]_{A}
  &= m^{4}\,\sigma_{2}^{6}\,(m-1)\bigl[(m^{3} - m_{2})\sigma_{2}^{2} + m_{2}\sigma_{4}\bigr]\\
[u^{(A, A)}_{\mathrm{init}}]_{U_{1}, U_{2}}
  &= m^{4}\,\sigma_{2}^{6}\,(m-1)\bigl[(m\,m_{1} + m - 2)\sigma_{2}^{2} + \sigma_{4}\bigr]\\
[u^{(A, A)}_{\mathrm{init}}]_{4K}
  &= m^{4}\,\sigma_{2}^{4}\bigl[
       (m^{3} - m\,m_{1} + 2m - m_{2} - 1)\sigma_{2}^{4}
       + (m\,m_{1} + m_{2} - 2)\sigma_{2}^{2}\sigma_{4} + \sigma_{4}^{2}\bigr]
\end{align*}

\paragraph{$u^{(A, U_{1})}_{\mathrm{init}}$.}
\begin{align*}
[u^{(A, U_{1})}_{\mathrm{init}}]_{A}
  &= m^{4}\,\sigma_{2}^{6}\,(m-1)\bigl[(m^{2} + m_{1} - 2)\sigma_{2}^{2} + \sigma_{4}\bigr]\\
[u^{(A, U_{1})}_{\mathrm{init}}]_{U_{1}}
  &= m^{4}\,\sigma_{2}^{6}\,(m-1)\bigl[(m^{2} + m\,m_{2} - 2m_{2})\sigma_{2}^{2} + m_{2}\sigma_{4}\bigr]\\
[u^{(A, U_{1})}_{\mathrm{init}}]_{U_{2}}
  &= m^{4}\,\sigma_{2}^{6}\,(m-1)\bigl[(2m + m_{1} - 3)\sigma_{2}^{2} + \sigma_{4}\bigr]\\
[u^{(A, U_{1})}_{\mathrm{init}}]_{4K}
  &= m^{4}\,\sigma_{2}^{4}\bigl[
       (2m^{2} + m\,m_{2} - m - m_{1} - 2m_{2} + 1)\sigma_{2}^{4}
       + (m + m_{1} + m_{2} - 3)\sigma_{2}^{2}\sigma_{4} + \sigma_{4}^{2}\bigr]
\end{align*}

\paragraph{$u^{(A, 4K)}_{\mathrm{init}}$.}
\begin{align*}
[u^{(A, 4K)}_{\mathrm{init}}]_{A}
  &= m^{4}\,\sigma_{2}^{4}\,(m-1)\bigl[
       (m^{3} - m^{2} + 2m_{1} - m_{2} - 1)\sigma_{2}^{4}
       + (m^{2} + m_{2} - 2)\sigma_{2}^{2}\sigma_{4} + \sigma_{4}^{2}\bigr]\\
[u^{(A, 4K)}_{\mathrm{init}}]_{U_{1}, U_{2}}
  &= m^{4}\,\sigma_{2}^{4}\,(m-1)\bigl[
       (m^{2} + m\,m_{1} + m\,m_{2} - 3m + m_{1} - 2m_{2} + 1)\sigma_{2}^{4}
       + (2m + m_{2} - 3)\sigma_{2}^{2}\sigma_{4} + \sigma_{4}^{2}\bigr]\\
[u^{(A, 4K)}_{\mathrm{init}}]_{4K}
  &= m^{4}\,\sigma_{2}^{2}\bigl[
       (m^{3} + m^{2} - m\,m_{1} + 2m\,m_{2} - 3m - 2m_{1} - 5m_{2} + 7)\sigma_{2}^{6}\\
  &\hspace{4em}
       + (m^{2} + m\,m_{1} + 2m_{1} + 3m_{2} - 7)\sigma_{2}^{4}\sigma_{4}
       + (m - 1)\sigma_{2}^{2}\sigma_{4}^{2} + \sigma_{4}^{3}\bigr]
\end{align*}

\paragraph{$u^{(U_{1}, A)}_{\mathrm{init}}$.}
\begin{align*}
[u^{(U_{1}, A)}_{\mathrm{init}}]_{A}
  &= m^{4}\,\sigma_{2}^{6}\,(m-1)\bigl[(m^{2} + m\,m_{2} - 2m_{2})\sigma_{2}^{2} + m_{2}\sigma_{4}\bigr]\\
[u^{(U_{1}, A)}_{\mathrm{init}}]_{U_{1}}
  &= m^{4}\,\sigma_{2}^{6}\,(m-1)\bigl[(m^{2} + m_{1} - 2)\sigma_{2}^{2} + \sigma_{4}\bigr]\\
[u^{(U_{1}, A)}_{\mathrm{init}}]_{U_{2}}
  &= m^{4}\,\sigma_{2}^{6}\,(m-1)\bigl[(2m + m_{1} - 3)\sigma_{2}^{2} + \sigma_{4}\bigr]\\
[u^{(U_{1}, A)}_{\mathrm{init}}]_{4K}
  &= m^{4}\,\sigma_{2}^{4}\bigl[
       (m^{2} + 2m\,m_{2} - m - m_{1} - 2m_{2} + 1)\sigma_{2}^{4}
       + (m + m_{1} + m_{2} - 3)\sigma_{2}^{2}\sigma_{4} + \sigma_{4}^{2}\bigr]
\end{align*}


\paragraph{$u^{(U_{1}, U_{1})}_{\mathrm{init}}$.}
\begin{align*}
[u^{(U_{1}, U_{1})}_{\mathrm{init}}]_{A}
  &= m^{4}\,\sigma_{2}^{6}\,(m-1)\bigl[(2m + m_{2} - 3)\sigma_{2}^{2} + \sigma_{4}\bigr]\\
[u^{(U_{1}, U_{1})}_{\mathrm{init}}]_{U_{1}}
  &= m^{4}\,\sigma_{2}^{6}\,(m-1)\bigl[(m^{3} - m_{1})\sigma_{2}^{2} + m_{1}\sigma_{4}\bigr]\\
[u^{(U_{1}, U_{1})}_{\mathrm{init}}]_{U_{2}}
  &= m^{4}\,\sigma_{2}^{6}\,(m-1)\bigl[(m + m_{1} + m_{2} - 3)\sigma_{2}^{2} + \sigma_{4}\bigr]\\
[u^{(U_{1}, U_{1})}_{\mathrm{init}}]_{4K}
  &= m^{4}\,\sigma_{2}^{4}\bigl[
       (m^{3} + m^{2} - m\,m_{2} - m_{1} - 2)\sigma_{2}^{4}
       + (m + m_{1} + m_{2} - 3)\sigma_{2}^{2}\sigma_{4} + \sigma_{4}^{2}\bigr]
\end{align*}

\paragraph{$u^{(U_{1}, U_{2})}_{\mathrm{init}}$.}
\begin{align*}
[u^{(U_{1}, U_{2})}_{\mathrm{init}}]_{A}
  &= m^{4}\,\sigma_{2}^{6}\,(m-1)\bigl[(2m + m_{2} - 3)\sigma_{2}^{2} + \sigma_{4}\bigr]\\
[u^{(U_{1}, U_{2})}_{\mathrm{init}}]_{U_{1}}
  &= m^{4}\,\sigma_{2}^{6}\,(m-1)\bigl[(m + m_{1} + m_{2} - 3)\sigma_{2}^{2} + \sigma_{4}\bigr]\\
[u^{(U_{1}, U_{2})}_{\mathrm{init}}]_{U_{2}}
  &= m^{4}\,\sigma_{2}^{6}\,(m-1)\bigl[(m + m_{1} + m_{2} - 3)\sigma_{2}^{2} + \sigma_{4}\bigr]\\
[u^{(U_{1}, U_{2})}_{\mathrm{init}}]_{4K}
  &= m^{4}\,\sigma_{2}^{4}\bigl[
       (2m^{2} - m + m_{1} + m_{2} - 4)\sigma_{2}^{4}
       + (m + m_{1} + m_{2} - 3)\sigma_{2}^{2}\sigma_{4} + \sigma_{4}^{2}\bigr]
\end{align*}

\paragraph{$u^{(U_{1}, 4K)}_{\mathrm{init}}$.}
\begin{align*}
[u^{(U_{1}, 4K)}_{\mathrm{init}}]_{A}
  &= m^{4}\,\sigma_{2}^{4}\,(m-1)\bigl[
       (m^{2} + m\,m_{1} + m\,m_{2} - 3m + m_{1} - 2m_{2} + 1)\sigma_{2}^{4}
       + (2m + m_{2} - 3)\sigma_{2}^{2}\sigma_{4} + \sigma_{4}^{2}\bigr]\\
[u^{(U_{1}, 4K)}_{\mathrm{init}}]_{U_{1}}
  &= m^{4}\,\sigma_{2}^{4}\,(m-1)\bigl[
       (m^{3} - m^{2} + 2m_{1} - m_{2} - 1)\sigma_{2}^{4}
       + (m^{2} + m_{2} - 2)\sigma_{2}^{2}\sigma_{4} + \sigma_{4}^{2}\bigr]\\
[u^{(U_{1}, 4K)}_{\mathrm{init}}]_{U_{2}}
  &= m^{4}\,\sigma_{2}^{4}\,(m-1)\bigl[
       (m^{2} + m\,m_{1} + m\,m_{2} - 3m + m_{1} - 2m_{2} + 1)\sigma_{2}^{4}
       + (2m + m_{2} - 3)\sigma_{2}^{2}\sigma_{4} + \sigma_{4}^{2}\bigr]\\
[u^{(U_{1}, 4K)}_{\mathrm{init}}]_{4K}
  &= m^{4}\,\sigma_{2}^{2}\bigl[
       (m^{3} + 2m^{2} - m\,m_{1} + m\,m_{2} - 6m + m_{1} - 2m_{2} + 4)\sigma_{2}^{6}\\
  &\hspace{4em}
       + (2m^{2} + m\,m_{1} - 4m - 2m_{1} - m_{2} + 4)\sigma_{2}^{4}\sigma_{4}
       + (m + m_{2} - 2)\sigma_{2}^{2}\sigma_{4}^{2} + \sigma_{4}^{3}\bigr]
\end{align*}

\paragraph{$u^{(4K, A)}_{\mathrm{init}}$.}
\begin{align*}
[u^{(4K, A)}_{\mathrm{init}}]_{A}
  &= m^{4}\,\sigma_{2}^{4}\,(m-1)\bigl[
       (m^{3} - m^{2} + 2m\,m_{2} - 2m_{2})\sigma_{2}^{4}
       + (m^{2} - m_{2})\sigma_{2}^{2}\sigma_{4} + m_{2}\sigma_{4}^{2}\bigr]\\
[u^{(4K, A)}_{\mathrm{init}}]_{U_{1}, U_{2}}
  &= m^{4}\,\sigma_{2}^{4}\,(m-1)\bigl[
       (m^{2} + m\,m_{1} - m_{1} - 1)\sigma_{2}^{4}
       + (m + m_{1} - 2)\sigma_{2}^{2}\sigma_{4} + \sigma_{4}^{2}\bigr]\\
[u^{(4K, A)}_{\mathrm{init}}]_{4K}
  &= m^{4}\,\sigma_{2}^{2}\bigl[
       (m^{3} + m^{2} - m\,m_{1} + 2m\,m_{2} - 6m + m_{1} - 2m_{2} + 4)\sigma_{2}^{6}\\
  &\hspace{4em}
       + (m^{2} + m\,m_{1} + 4m - 2m_{1} - m_{2} - 3)\sigma_{2}^{4}\sigma_{4}
       + (m_{1} + m_{2} - 2)\sigma_{2}^{2}\sigma_{4}^{2} + \sigma_{4}^{3}\bigr]
\end{align*}

\paragraph{$u^{(4K, U_{1})}_{\mathrm{init}}$.}
\begin{align*}
[u^{(4K, U_{1})}_{\mathrm{init}}]_{A}
  &= m^{4}\,\sigma_{2}^{4}\,(m-1)\bigl[
       (m^{2} + m\,m_{1} - m_{1} - 1)\sigma_{2}^{4}
       + (m + m_{1} - 2)\sigma_{2}^{2}\sigma_{4} + \sigma_{4}^{2}\bigr]\\
[u^{(4K, U_{1})}_{\mathrm{init}}]_{U_{1}}
  &= m^{4}\,\sigma_{2}^{4}\,(m-1)\bigl[
       (m^{3} - m\,m_{1} + 2m\,m_{2} - 2m_{2})\sigma_{2}^{4}
       + (m\,m_{1} - m_{2})\sigma_{2}^{2}\sigma_{4} + m_{2}\sigma_{4}^{2}\bigr]\\
[u^{(4K, U_{1})}_{\mathrm{init}}]_{U_{2}}
  &= m^{4}\,\sigma_{2}^{4}\,(m-1)\bigl[
       (m^{2} + m\,m_{1} - m_{1} - 1)\sigma_{2}^{4}
       + (m + m_{1} - 2)\sigma_{2}^{2}\sigma_{4} + \sigma_{4}^{2}\bigr]\\
[u^{(4K, U_{1})}_{\mathrm{init}}]_{4K}
  &= m^{4}\,\sigma_{2}^{2}\bigl[
       (m^{3} + m^{2} - m\,m_{1} + 2m\,m_{2} - 6m + m_{1} - 2m_{2} + 4)\sigma_{2}^{6}\\
  &\hspace{4em}
       + (m^{2} + m\,m_{1} + 4m - 2m_{1} - m_{2} - 3)\sigma_{2}^{4}\sigma_{4}
       + (m_{1} + m_{2} - 2)\sigma_{2}^{2}\sigma_{4}^{2} + \sigma_{4}^{3}\bigr]
\end{align*}

\paragraph{$u^{(4K, 4K)}_{\mathrm{init}}$.}
\begin{align*}
[u^{(4K, 4K)}_{\mathrm{init}}]_{A, U_{1}, U_{2}}
  &= m^{4}\,\sigma_{2}^{2}\,(m-1)\bigl[
       (m^{3} - 2m^{2} + 2m\,m_{1} + 2m\,m_{2} - 3m - 2m_{1} - 2m_{2} + 4)\sigma_{2}^{6}\\
  &\hspace{4em}
       + (2m^{2} + 2m_{1} - m_{2} - 3)\sigma_{2}^{4}\sigma_{4}
       + (m + m_{2} - 2)\sigma_{2}^{2}\sigma_{4}^{2} + \sigma_{4}^{3}\bigr]\\
[u^{(4K, 4K)}_{\mathrm{init}}]_{4K}
  &= m^{4}\bigl[
       3(m^{3} - 2m^{2} - m\,m_{1} + 2m\,m_{2} + m_{1} - 2m_{2} + 1)\sigma_{2}^{8}\\
  &\hspace{2em}
       + 3(2m^{2} + m\,m_{1} - 4m - 2m_{1} - m_{2} + 4)\sigma_{2}^{6}\sigma_{4}
       + 3(2m + m_{1} + m_{2} - 4)\sigma_{2}^{4}\sigma_{4}^{2} + \sigma_{4}^{4}\bigr]
\end{align*}
\paragraph{End vector $u^{P_{N}}_{\mathrm{end}}$.}

We close the chain. At this point, the contributions of sites $1, \dots, N-1$ 
are encoded in a vector indexed by the pattern of the bond $z_{N-2}$, obtained 
by contracting $u^{(P_{1}, P_{2})}_{\mathrm{init}}$ with the bulk transfer 
matrices of Theorem~\ref{thm:W-bulk}. The remaining task is to absorb 
the boundary site $N$, consisting of leaf $X^{(N)}$ and the last gadget 
$g^{(N-1)}$, into a vector $u^{P_{N}}_{\mathrm{end}}$ that contracts with 
the chain state to produce a scalar.

Unlike interior gadgets, $g^{(N-1)}$ has no outgoing bond $z_{N-1}$: 
its fourth index is the output index $n$ of the sketching tensor 
network~\eqref{eq:g_def}. The boundary trace structure of the original 
sum~\eqref{eq:S_4_copies} pairs the four replicas of $n$ as 
$n^{(a)} = n^{(b)}$ and $n^{(c)} = n^{(d)}$, so the pattern of $n$ is 
restricted to $\{A, 4K\}$ and is summed over rather than fixed.

\begin{theorem}[End vector]\label{thm:u-end}
Fix $N \ge 3$ and a pattern $P_{N} \in \mathcal{P}$ for the outer index of the boundary leaf $X^{(N)}$. The joint contribution of leaf $X^{(N)}$ and gadget $g^{(N-1)}$ to the sum~\eqref{eq:S_4_copies}, conditional on 
the incoming bond pattern $Z_{\mathrm{in}}$ of $z_{N-2}$, is the entry $\bigl[u^{P_{N}}_{\mathrm{end}}\bigr]_{Z_{\mathrm{in}}}$ of a vector $u^{P_{N}}_{\mathrm{end}} \in \mathbb{R}^{4}$ indexed in the order $(A, U_{1}, U_{2}, 4K)$, with
\begin{equation}\label{eq:u-end}
\bigl[u^{P_{N}}_{\mathrm{end}}\bigr]_{Z_{\mathrm{in}}}
\;=\;
\sum_{\mathcal{N} \in \{A, 4K\}} N_{\mathcal{N}}(m)
\!\!\sum_{K\,\in\, \mathcal{K}(P_{N})}\!\! \mu^{X}_{P_{N}, K}
\!\!\!\!\sum_{\substack{K_{1}, K_{2} \,\in\, \mathcal{P} \\ K_{1} \wedge K_{2} \,=\, K}}\!\!\!\!
N_{K_{1}}(m_{1})\, N_{K_{2}}(m_{2})\,
\mu^{g}_{\,Z_{\mathrm{in}} \wedge K_{1},\, K_{2} \wedge \mathcal{N}}.
\end{equation}
\end{theorem}

\begin{proof}
\textbf{Setting up the model.} 
We close the chain by enumerating site $N$, which consists of $X^{(N)}$ 
and $g^{(N-1)}$. We remember the indices of $g^{(N-1)}_{z_{N-2} k^{(1)}_{N} k^{(2)}_{N} n}$, 
eq.~\eqref{eq:g_def}, and the leaf $X^{(N)}_{k_{N}, i_{N}}$, eq.~\eqref{eq:S_4_copies}. 
Note that, unlike interior gadgets, the right outer index of $g^{(N-1)}$ 
is the output index $n$ rather than an outgoing bond.

The restrictions are the same as before, with an added restriction the arises from the inner product structure of eq. (\ref{eq:S_dot_prod_def}). We get that $$n^{(a)}=n^{(b)}\land n^{(c)}=n^{(d)}$$
\smallskip\noindent\textbf{Gadget layer.}
The right half $G^{(N-1, 2)}$ now carries the indices $(k^{(2)}_{N}, n)$ 
instead of $(k^{(2)}_{p+1}, z_{p})$, so the right half pattern is 
$R = K_{2} \wedge \mathcal{N}$ where $\mathcal{N}$ is the pattern of $n$. 
The left half is unchanged, $L = Z_{\mathrm{in}} \wedge K_{1}$.

\smallskip\noindent\textbf{Leaf layer.}
Identical to Theorem~\ref{thm:W-bulk}: the leaf $X^{(N)}$ contributes 
$\mu^{X}_{P_{N}, K}$ where $K:= K_{1} \wedge K_{2}$, non-vanishing only 
when $K \in \mathcal{K}(P_{N})$.

\smallskip\noindent\textbf{Output-index layer.}
In place of the outgoing-bond layer, the output index $n$ is summed 
rather than fixed: for each pattern $\mathcal{N} \in \{A, 4K\}$ admitted 
by the boundary trace, there are $N_{\mathcal{N}}(m)$ value-tuples 
by Lemma~\ref{lem:N-count}, and we sum over both admissible patterns.

\medskip
Collecting the three layers yields~\eqref{eq:u-end}.
\end{proof}
\subsubsection{ Explicit End vectors $u_{end}$}
\paragraph{$u^{A}_{\mathrm{end}}$.}
\begin{align*}
[u^{A}_{\mathrm{end}}]_{A}
  &= m^{3}\,\sigma_{2}^{4}\bigl[(m^{3} - m_{1})\sigma_{2}^{2} + m_{1}\sigma_{4}\bigr]\\
[u^{A}_{\mathrm{end}}]_{U_{1}, U_{2}}
  &= m^{3}\,\sigma_{2}^{4}\bigl[(m\,m_{2} + m - 2)\sigma_{2}^{2} + \sigma_{4}\bigr]\\
[u^{A}_{\mathrm{end}}]_{4K}
  &= m^{3}\,\sigma_{2}^{2}\bigl[
       (m^{3} - m\,m_{2} + 2m - m_{1} - 1)\sigma_{2}^{4}
       + (m\,m_{2} + m_{1} - 2)\sigma_{2}^{2}\sigma_{4} + \sigma_{4}^{2}\bigr]
\end{align*}

\paragraph{$u^{U_{1}}_{\mathrm{end}}$.}
\begin{align*}
[u^{U_{1}}_{\mathrm{end}}]_{A}
  &= m^{3}\,\sigma_{2}^{4}\bigl[(m^{2} + m_{2} - 2)\sigma_{2}^{2} + \sigma_{4}\bigr]\\
[u^{U_{1}}_{\mathrm{end}}]_{U_{1}}
  &= m^{3}\,\sigma_{2}^{4}\bigl[(m^{2} + m\,m_{1} - 2m_{1})\sigma_{2}^{2} + m_{1}\sigma_{4}\bigr]\\
[u^{U_{1}}_{\mathrm{end}}]_{U_{2}}
  &= m^{3}\,\sigma_{2}^{4}\bigl[(2m + m_{2} - 3)\sigma_{2}^{2} + \sigma_{4}\bigr]\\
[u^{U_{1}}_{\mathrm{end}}]_{4K}
  &= m^{3}\,\sigma_{2}^{2}\bigl[
       (2m^{2} + m\,m_{1} - m - 2m_{1} - m_{2} + 1)\sigma_{2}^{4}
       + (m + m_{1} + m_{2} - 3)\sigma_{2}^{2}\sigma_{4} + \sigma_{4}^{2}\bigr]
\end{align*}

\paragraph{$u^{4K}_{\mathrm{end}}$.}
\begin{align*}
[u^{4K}_{\mathrm{end}}]_{A}
  &= m^{3}\,\sigma_{2}^{2}\bigl[
       (m^{3} - m^{2} - m_{1} + 2m_{2} - 1)\sigma_{2}^{4}
       + (m^{2} + m_{1} - 2)\sigma_{2}^{2}\sigma_{4} + \sigma_{4}^{2}\bigr]\\
[u^{4K}_{\mathrm{end}}]_{U_{1}, U_{2}}
  &= m^{3}\,\sigma_{2}^{2}\bigl[
       (m^{2} + m\,m_{1} + m\,m_{2} - 3m - 2m_{1} + m_{2} + 1)\sigma_{2}^{4}
       + (2m + m_{1} - 3)\sigma_{2}^{2}\sigma_{4} + \sigma_{4}^{2}\bigr]\\
[u^{4K}_{\mathrm{end}}]_{4K}
  &= m^{3}\bigl[
       (m^{3} + m^{2} + 2m\,m_{1} - m\,m_{2} - 3m - 5m_{1} - 2m_{2} + 7)\sigma_{2}^{6}\\
  &\hspace{2em}
       + (m^{2} + m\,m_{2} + 3m_{1} + 2m_{2} - 7)\sigma_{2}^{4}\sigma_{4}
       + (m - 1)\sigma_{2}^{2}\sigma_{4}^{2} + \sigma_{4}^{3}\bigr]
\end{align*}
\subsection{Glueing the Theorems Together}
\begin{proof}[Proof of Theorem~\ref{thm:per-entry-M}]
The index quadruple $(i^{a},i^{b},i^{c},i^{d})$ induces a chain of leaf patterns $(P_{1},\dots,P_{N})\in\mathcal{P}^{N}$ (Def.~\ref{def:Chain}). Applying Theorems~\ref{thm:u-init},~\ref{thm:W-bulk}, and~\ref{thm:u-end} to the corresponding sites yields~\eqref{eq:chain-value}.
\end{proof}
\pagebreak
\section{Bounding $Tr(MB)$}

\textbf{Lemma} 
A matrix with the structure $\mathbb{E}[S^{T}S\otimes S^{T}S]$ for some random matrix $S\in \R^{m\times n}$   is a PSD matrix.
\textbf{Proof}
Given some $x\in\R^{n}$ will calculate $x^TMx$ and prove its non-negative.
\begin{align}
    x^TMx&=\sum_{i^{(a)},i^{(c)}=1}^d\sum_{i^{(b)},i^{(d)}=1}^dx_{i^{(a)},i^{(c)}}M_{i^{(a)},i^{(b)},i^{(c)},i^{(d)}}x_{i^{(b)},i^{(d)}}\\
    &=\sum_{i^{(a)},i^{(c)}=1}^d\sum_{i^{(b)},i^{(d)}=1}^dx_{i^{(a)},i^{(c)}}\mathbb{E}[\langle S_{i^{(a)}},S_{i^{(b)}}\rangle\cdot\langle S_{i^{(c)}},S_{i^{(d)}}\rangle]x_{i^{(b)},i^{(d)}}\\
    &=\mathbb{E}[\sum_{i^{(a)},i^{(c)}=1}^d\sum_{i^{(b)},i^{(d)}=1}^dx_{i^{(a)},i^{(c)}}\langle S_{i^{(a)}},S_{i^{(b)}}\rangle\cdot\langle S_{i^{(c)}},S_{i^{(d)}}\rangle x_{i^{(b)},i^{(d)}}]\\
    &=\mathbb{E}[\sum_{i^{(a)},i^{(c)}=1}^d\sum_{i^{(b)},i^{(d)}=1}^dx_{i^{(a)},i^{(c)}}\sum_{n^{(1)}=1}^m S_{i^{(a)}n^{(1)}}S_{i^{(b)}n^{(1)}}\cdot\sum_{n^{(2)}=1}^m S_{i^{(c)}n^{(2)}},S_{i^{(d)}} x_{i^{(b)},i^{(d)}}]\\
&=\mathbb{E}[\sum_{n^{(1)}=1}^m\sum_{n^{(2)}=1}^m\sum_{i^{(a)},i^{(c)}=1}^d\sum_{i^{(b)},i^{(d)}=1}^dx_{i^{(a)},i^{(c)}} S_{i^{(a)}n^{(1)}},S_{i^{(b)}n^{(1)}} S_{i^{(c)}n^{(2)}}S_{i^{(d)}} x_{i^{(b)},i^{(d)}}]\\
&=\mathbb{E}[\sum_{n^{(1)}=1}^m\sum_{n^{(2)}=1}^m(\sum_{i^{(a)},i^{(c)}=1}^dx_{i^{(a)},i^{(c)}} S_{i^{(a)}n^{(1)}}S_{i^{(c)}n^{(2)}})(\sum_{i^{(b)},i^{(d)}=1}^dS_{i^{(b)}n^{(1)}}\cdot S_{i^{(d)}} x_{i^{(b)},i^{(d)}})]\\
&=\mathbb{E}[\sum_{n^{(1)}=1}^m\sum_{n^{(2)}=1}^mY_{n^1,n^2}^2]\\
&\ge0
\end{align}
Where $Y=\sum_{i^{(a)},i^{(c)}=1}^dx_{i^{(a)},i^{(c)}} S_{i^{(a)}n^{(1)}}S_{i^{(c)}n^{(2)}}=\sum_{i^{(b)},i^{(d)}=1}^dx_{i^{(b)},i^{(d)}} S_{i^{(b)}n^{(1)}}S_{i^{(d)}n^{(2)}}$.
\begin{corollary}\label{lem:M-psd}
With this lemma, we get that $M$ is a PSD and we also know that $B=A\otimes A$ is a PSD, see eq(\ref{eq:M_def}).
Meaning we can use the bounds:
\begin{itemize}
    \item $Tr(MB)\le Tr(M)Tr(B)$
    \item $Tr(MB)\le \|M\|_\infty Tr(B)$ 
\end{itemize}
\end{corollary}

\subsection{Bound Via $\mathrm{Tr}(M)$}
\begin{theorem}[Closed form for $\mathrm{Tr}(M)$]
Given $M$ like in eq.(\ref{eq:M_def}) with $N\ge 3$ and using the vectors and matrices from theorem (\ref{thm:per-entry-M}).
Define
\begin{align}
\mathbf{u}^{\mathrm{Tr}}_{\mathrm{init}}
&:=\;d^{2}\,u^{(4K,4K)}_{\mathrm{init}}
\;+\;d^{2}(d-1)\bigl(u^{(4K,A)}_{\mathrm{init}}+u^{(A,4K)}_{\mathrm{init}}\bigr)
\;+\;d^{2}(d-1)^{2}\,u^{(A,A)}_{\mathrm{init}},\\[2pt]
\mathbf{u}^{\mathrm{Tr}}_{\mathrm{end}}
&:=\;d\,u^{4K}_{\mathrm{end}}\;+\;d(d-1)\,u^{A}_{\mathrm{end}},\\[2pt]
\mathbf{W}^{\mathrm{Tr}}
&:=\;d\,W^{4K}\;+\;d(d-1)\,W^{A}.
\end{align}
Then
\begin{align}
\mathrm{Tr}(M)\;=\;\bigl\langle\,\mathbf{u}^{\mathrm{Tr}}_{\mathrm{end}},\;
(\mathbf{W}^{\mathrm{Tr}})^{N-3}\,\mathbf{u}^{\mathrm{Tr}}_{\mathrm{init}}\,\bigr\rangle. \label{def:tr_formula}
\end{align}


For $N=3$, $(\mathbf{W}^{\mathrm{Tr}})^{0}=I$ and
$\mathrm{Tr}(M)=\langle\mathbf{u}^{\mathrm{Tr}}_{\mathrm{end}},\mathbf{u}^{\mathrm{Tr}}_{\mathrm{init}}\rangle$.
\end{theorem}
\begin{proof}
By definition,
\begin{align}
{Tr}(M)
=\sum_{\boldsymbol{i}^{(a)},\boldsymbol{i}^{(c)}\in[d]^{N}}
M_{\boldsymbol{i}^{(a)},\boldsymbol{i}^{(a)},\boldsymbol{i}^{(c)},\boldsymbol{i}^{(c)}}\label{def:tr_sum}
\end{align}


We note that because of the trace condition on the indices, $i^{(a)}=i^{(b)}\land i^{(c)}=i^{(d)}$, we get that the pattern chain for each element, eq (\ref{def:Chain}),  contains only patterns $\{4K,A\}$. i.e. each element is of the form $\mathcal{P}=(P_{1},\ldots,P_{N})\in\{4K,A\}^{N}$.\\
We get from theorem (\ref{eq:per-entry-M}) that each element on the main diagonal is of the form
\begin{align}
M_{\boldsymbol{i}^{(a)},\boldsymbol{i}^{(a)},\boldsymbol{i}^{(c)},\boldsymbol{i}^{(c)}}=\langle u_{end}^{P_N},\prod_{j=3}^{N-1}W^{p_j}u_{start}^{(P_1,P_2)}\rangle
\end{align}
With this in mind we can change the sum from summing over the indices to summing over the pattern chains where we count how many times each pattern chain occurs.

\paragraph{Count Per main diagonal pattern chain occurrences}
For a fixed pattern sequence $\boldsymbol{P}\in\{4K,A\}^{N}$, count the index pairs $(\boldsymbol{i}^{(a)},\boldsymbol{i}^{(c)})\in[d]^{N}\times[d]^{N}$ that realise it.
Each site $p=1,..,N$ is independent, if $P_p=4K$ then $i^{(a)}_p=i^{(c)}_p$, giving $d$ choices.
If $P_p=A$ then $i^{(a)}_p\ne i^{(c)}_p$, giving $d(d-1)$ choices. 
Regrouping \eqref{def:tr_sum} by pattern chain, with multiplicity $\prod_{p=1}^{N} c_{P_p}$ where $c_{4K}:=d,\ c_A:=d(d-1)$, and using linearity of $\langle\cdot,\cdot\rangle$ to absorb each site sum into its factor:
\begin{align}
\mathrm{Tr}(M)
&= \sum_{\mathcal{P}\in\{4K,A\}^{N}}\!\Bigl(\prod_{p=1}^{N} c_{P_p}\Bigr)
\bigl\langle u_{end}^{P_N},\, W^{P_{N-1}}\!\cdots W^{P_3}\, u_{start}^{(P_1,P_2)}\bigr\rangle \notag\\
&= \Bigl\langle\, \sum_{P_N} c_{P_N} u_{end}^{P_N},\;
\Bigl(\!\!\sum_{P_3,\ldots,P_{N-1}}\!\prod_{j=3}^{N-1}c_{P_j}\,W^{P_{N-1}}\!\cdots W^{P_3}\Bigr)
\sum_{P_1,P_2} c_{P_1}c_{P_2}\, u_{start}^{(P_1,P_2)} \Bigr\rangle,
\end{align}
The first and last sums equal $\mathbf{u}_{end}^{\mathrm{Tr}}$ and $\mathbf{u}_{start}^{\mathrm{Tr}}$ by definition. For the middle, the non-commutative binomial expansion of a product of $N-3$ identical sums gives
\begin{align}
(\mathbf{W}^{\mathrm{Tr}})^{N-3}
\;=\;\Bigl(\sum_{P\in\{4K,A\}}c_P W^P\Bigr)^{\!N-3}
\;=\!\!\sum_{(P_3,\ldots,P_{N-1})\in\{4K,A\}^{N-3}}\!\!\prod_{j=3}^{N-1}c_{P_j}\;W^{P_{N-1}}\!\cdots W^{P_3},
\end{align}
Substituting into \eqref{eq:tr-pulled-in},
\begin{align}
\mathrm{Tr}(M)=\bigl\langle \mathbf{u}_{end}^{\mathrm{Tr}},\,(\mathbf{W}^{\mathrm{Tr}})^{N-3}\,\mathbf{u}_{start}^{\mathrm{Tr}}\bigr\rangle.
\end{align}

\end{proof}

\subsection{Bound via $\|M\|_{\infty }$}

Since $M \text{ and } B$ are both PSD's, see (\ref{lem:M-psd}), we can use the following bound
\begin{equation}
\mathrm{Tr}(MB)\le\;\|M\|_\infty \,\mathrm{Tr}(B),
\end{equation}
In order to find $\|M\|_\infty$  we do the following:
1) Use Greshgorin's theorem to bound the eigenvalues with maximum row sum
2) Proof that the maximum row sum is indeed an eigenvalue.
With these $2$  step we get the largest eigenvalues and hence the operator norm for $M$ with the following theorem:

\begin{theorem}[Closed form for $\|M\|_\infty$]\label{thm:opnorm-closed}
Let $M$ be the matrix of eq.~\eqref{eq:M_def} with $N\ge 3, d,m\ge 2$,  whose entries are
given by the per entry Theorem~\ref{thm:per-entry-M}, and let $\{W^{P},\,u^{(P_1,P_2)}_{\mathrm{init}},\,u^{P_N}_{\mathrm{end}}\}_{P\in\mathcal{P}}$ be the associated transfer matrices and boundary vectors. 
Define 
\begin{align}
\mathbf{u}^{\mathrm{op}}_{\mathrm{init}}
&:=\;u^{(4K,4K)}_{\mathrm{init}}
\;+\;(d-1)\bigl(u^{(4K,U_1)}_{\mathrm{init}}+u^{(U_1,4K)}_{\mathrm{init}}\bigr)
\;+\;(d-1)^{2}\,u^{(U_1,U_1)}_{\mathrm{init}},\\[2pt]
\mathbf{u}^{\mathrm{op}}_{\mathrm{end}}
&:=\;u^{4K}_{\mathrm{end}}\;+\;(d-1)\,u^{U_1}_{\mathrm{end}},\\[2pt]
\mathbf{W}^{\mathrm{op}}
&:=\;W^{4K}\;+\;(d-1)\,W^{U_1}.
\end{align}
Then for $M$
\begin{equation}\label{eq:opnorm_formula}
\;\|M\|_{\infty}
=\;\bigl\langle\,\mathbf{u}^{\mathrm{op}}_{\mathrm{end}},\;
(\mathbf{W}^{\mathrm{op}})^{N-3}\,\mathbf{u}^{\mathrm{op}}_{\mathrm{init}}\,\bigr\rangle .
\end{equation}
For $N=3$, $(\mathbf{W}^{\mathrm{op}})^{0}=I$ and
$\|M\|_\infty=\langle\mathbf{u}^{\mathrm{op}}_{\mathrm{end}},\mathbf{u}^{\mathrm{op}}_{\mathrm{init}}\rangle$.
\end{theorem}
\subsubsection{Evaluate the per row Sum}
\paragraph{Row classification by $E/D$.}
\begin{definition}[$E/D$ pattern]\label{def:E_D_sites}
Fix a row $(i^{(a)}, i^{(c)})\in\R^{d^N}\times\R^{d^N}$ of $M$. For each site $l\in[N]$, set
\begin{equation}
s_l \;:=\; \begin{cases} E, & i_l^{(a)} = i_l^{(c)},\\ D, & i_l^{(a)} \neq i_l^{(c)}. \end{cases}
\end{equation}
$E$  for Equal and $D$  for Different.
The vector $\mathbf{s} = (s_1,\dots,s_N)\in\{E,D\}^N$ is the \emph{$E/D$-pattern} of the row.
\end{definition}

\begin{lemma}[Per-site compatibility]\label{lem:per_site_compa}
Fix a row with $E/D$-pattern $\mathbf{s}$ and a column $(i^{(b)}, i^{(d)})$.
The corresponding entry of $M$ is nonzero if and only if at every site $l$ the joint $(i_l^{(a)}, i_l^{(b)}, i_l^{(c)}, i_l^{(d)})$ has pattern $P_l\in\mathcal{P}$ where:
\begin{itemize}
\item if $s_l = E$, then $P_l \in \{4K, U_1\}$, with $1$ column choice for $P_l = 4K$ and $d-1$ choices for $P_l = U_1$;
\item if $s_l = D$, then $P_l \in \{A, U_2\}$, each with $1$ column choice.
\end{itemize}
\end{lemma}
\begin{proof}
We know that for each element in $M$  if its pattern chain, see def (\ref{def:Chain}), doesn't consist of only pattern in $\mathcal{P}$ , see def (\ref{def:patterns}), the element is $0$.
We look at some $E/D$-pattern for a row $(i^{(a)}.i^{(b)})$, we look at a site $l$  for each element and conclude what pattern that site must have.
\smallskip\noindent\textbf{Case $s_l=E$ ($i_l^{(a)}=i_l^{(c)}$).}
Because $i^{(a)}_l=i^{(c)}$  then $P_l\not\in\{A,U_2\}$ thus at an $E$-site only $P_l\in\{4K,U_1\}$ can give a nonzero moment. Counting the values $(i_l^{(b)},i_l^{(d)})$ consistent with each:
\begin{itemize}
\item $P_l=4K$ forces $i_l^{(b)}=i_l^{(d)}=i_l^{(a)}$, a single choice.
\item $P_l=U_1$ forces $i_l^{(a)}=i_l^{(c)}$ and $i_l^{(b)}=i_l^{(d)}$ with $i_l^{(b)}\neq i_l^{(a)}$, i.e.\ $i_l^{(b)}=i_l^{(d)}$ giving $d-1$ choices.
\end{itemize}

\smallskip\noindent\textbf{Case $s_l=D$ ($i_l^{(a)}\neq i_l^{(c)}$).}
Symmetric, exchanging the roles of the ``equal'' and ``different'' constraints. Now $4K$ and $U_1$ both contradict $i_l^{(a)}\neq i_l^{(c)}$ and contribute $0$, while $P_l\in\{A,U_2\}$ are admissible: $A$ forces $i_l^{(b)}=i_l^{(a)}$, $i_l^{(d)}=i_l^{(c)}$ and $U_2$ forces $i_l^{(d)}=i_l^{(a)}$, $i_l^{(b)}=i_l^{(c)}$, each pinning $(i_l^{(b)},i_l^{(d)})$ to a single value. Hence one column choice apiece.
\end{proof}

\begin{definition}
    For an $E/D$-pattern $\mathbf{s}$, see (\ref{def:E_D_sites}), define
\begin{equation}\label{eq:Rdef}
R(\mathbf{s}) \;:=\; \sum_{j} M_{ij}
\end{equation}
denote the sum of any row $i$ of $M$ whose pattern is $\mathbf{s}$.
The following theorem will show that $R(s)$  is well defined.
\end{definition}

\begin{definition}[Aggregated operators]\label{def:row_sum_operators}
Define
\begin{equation}
W^{E} := W^{4K} + (d-1)\,W^{U_1},
\qquad
W^{D} := W^{A} + W^{U_2},
\end{equation}
\begin{equation}
u_\mathrm{end}^{E} := u_\mathrm{end}^{4K} + (d-1)\,u_\mathrm{end}^{U_1},
\qquad
u_\mathrm{end}^{D} := u_\mathrm{end}^{A} + u_\mathrm{end}^{U_2},
\end{equation}
\begin{align}
u_\mathrm{init}^{(E,E)}
&:= u_\mathrm{init}^{(4K,4K)}
   + (d-1)\bigl(u_\mathrm{init}^{(4K,U_1)} + u_\mathrm{init}^{(U_1,4K)}\bigr)
   + (d-1)^2\,u_\mathrm{init}^{(U_1,U_1)},\\[2pt]
u_\mathrm{init}^{(E,D)}
&:= u_\mathrm{init}^{(4K,A)} + u_\mathrm{init}^{(4K,U_2)}
   + (d-1)\bigl(u_\mathrm{init}^{(U_1,A)} + u_\mathrm{init}^{(U_1,U_2)}\bigr),\\[2pt]
u_\mathrm{init}^{(D,E)}
&:= u_\mathrm{init}^{(A,4K)} + u_\mathrm{init}^{(U_2,4K)}
   + (d-1)\bigl(u_\mathrm{init}^{(A,U_1)} + u_\mathrm{init}^{(U_2,U_1)}\bigr),\\[2pt]
u_\mathrm{init}^{(D,D)}
&:= u_\mathrm{init}^{(A,A)} + u_\mathrm{init}^{(A,U_2)}
   + u_\mathrm{init}^{(U_2,A)} + u_\mathrm{init}^{(U_2,U_2)}.
\end{align}

\end{definition}


\begin{theorem}[Row-sum decomposition]\label{thm:row-decomp}
Fix an $E/D$-pattern $\mathbf{s}\in\{E,D\}^N$, def (\ref{def:E_D_sites}), and call a chain $\mathbf{P}\in\mathcal{P}^N$ \emph{compatible} with $\mathbf{s}$ if every $P_l$ is admissible at site $l$ in the sense of Lemma~\ref{lem:per_site_compa}. Write
\[
\mathcal{C}(\mathbf{s}) \;:=\; \{\mathbf{P}\in\mathcal{P}^N : \mathbf{P}\text{ compatible with }\mathbf{s}\},
\qquad
n_{\mathbf{P}} \;:=\; (d-1)^{\,\#\{l\,:\,P_l = U_1\}},
\]
for the set of \emph{compatible chains} and the \emph{chain multiplicity}. Then $R(\mathbf{s})$ in \eqref{eq:Rdef} is independent of the chosen row and
\begin{equation}\label{eq:row-decomp}
R(\mathbf{s}) \;=\; \sum_{\mathbf{P} \in \mathcal{C}(\mathbf{s})} n_{\mathbf{P}}\, V(\mathbf{P}).
\end{equation}
With the aggregated transfer matrices and boundary vectors of Definition~(\ref{def:row_sum_operators})
\begin{equation}\label{eq:row-transfer}
R(\mathbf{s}) \;=\; \bigl\langle u_\mathrm{end}^{(s_N)},\; W^{(s_{N-1})}\cdots W^{(s_3)}\, u_\mathrm{init}^{(s_1,s_2)}\bigr\rangle.
\end{equation}
\end{theorem}

\begin{proof}
Fix a row $i=(i^{(a)},i^{(c)})$ with pattern $\mathbf{s}$ and sum over columns $j=(i^{(b)},i^{(d)})$. Each index $(i^{(a)},i^{(b)},i^{(c)},i^{(d)})$ induces a chain $\mathbf{P}\in\mathcal{P}^N$ (Def.~\ref{def:Chain}), and by Theorem~\ref{thm:per-entry-M} the entry depends only on this chain, $V(\mathbf{P})$, with $V$ given by \eqref{eq:chain-value}; in particular it is independent of the underlying index values.

By Lemma~\ref{lem:per_site_compa}, $V(\mathbf{P})\neq 0$ iff $\mathbf{P}\in\mathcal{C}(\mathbf{s})$, and for each site we have $(d-1)$  options if $\mathbf{P}_l$  os of type $U_1$  and $1$  option otherwise. For a fixed admissible chain $\mathbf{P}$ the column entries are independent across each site, so the number of possible columns realizing $\mathbf{P}$ is
\[
\prod_{l=1}^N \#\{(i_l^{(b)},i_l^{(d)}) : \text{pattern }P_l=U_1\}
\;=\; (d-1)^{\#\{l:P_l=U_1\}} \;=\; n_{\mathbf{P}},
\]
since each $U_1$-site admits $d-1$ choices and every other admissible pattern admits exactly one. Changing summation from columns to all possible $\mathbf{P}$  gives.
\[
R(\mathbf{s})=\sum_j M_{ij}=\sum_{\mathbf{P}\in\mathcal{C}(\mathbf{s})} n_{\mathbf{P}}\,V(\mathbf{P}),
\]
which proves \eqref{eq:row-decomp}; the right-hand side depends only on $\mathbf{s}$, establishing row-independence.
Denote for convinience
\[
w_{P_l}(s_l)=
\begin{cases}
d-1, & (s_l,P_l)=(E,U_1),\\
1, & \text{otherwise},
\end{cases}
\qquad\Longrightarrow\qquad
n_{\mathbf{P}}=\prod_{l=1}^{N} w_{P_l}(s_l).
\]
Inserting this and the factorization \eqref{eq:chain-value} of $V(\mathbf{P})$ from Theorem~\ref{thm:per-entry-M} into \eqref{eq:row-decomp},
\[
R(\mathbf{s})
=\!\!\sum_{P_1\in\mathcal{P}(s_1)}\!\!\cdots\!\!\sum_{P_N\in\mathcal{P}(s_N)}\!\!\Bigl(\textstyle\prod_{l=1}^{N}w_{P_l}(s_l)\Bigr)\,\bigl(u_\mathrm{end}^{P_N}\bigr)^{\!\top} W^{P_{N-1}}\!\cdots W^{P_3}\,u_\mathrm{init}^{(P_1,P_2)}.
\]
Each $P_l$ occurs in exactly one factor, so the sum factorizes site by site:
\[
R(\mathbf{s})
=\Bigl(\!\sum_{P_N}\! w_{P_N}(s_N)\,u_\mathrm{end}^{P_N}\Bigr)^{\!\top}
\Bigl(\!\prod_{l=3}^{N-1}\sum_{P_l}\! w_{P_l}(s_l)\,W^{P_l}\Bigr)
\Bigl(\!\sum_{P_1,P_2}\! w_{P_1}(s_1)w_{P_2}(s_2)\,u_\mathrm{init}^{(P_1,P_2)}\Bigr).
\]
Plugging in definition~\ref{def:row_sum_operators} the three brackets are $u_\mathrm{end}^{(s_N)}$, $\prod_{l=3}^{N-1}W^{(s_l)}$, and $u_\mathrm{init}^{(s_1,s_2)}$, giving \eqref{eq:row-transfer}.
\end{proof}

\begin{corollary}\label{cor:EN-rowsum}
For each row of $M$ with indices $(i,i)$, for some $i\in d^N$ its sum equals $R(E^N)$.
\end{corollary}

\subsubsection{Find the maximal row sum}
Having established the per-row formula \eqref{eq:row-transfer} for $R(\mathbf{s})$, we identify the pattern $\mathbf{s}$ maximizing it.

\begin{lemma}[Entrywise dominance $W^E \succeq W^D$]\label{lem:transition_matrix_dominance}
For $\sigma_4 \geq \sigma_2^2$ and $d\geq 2$,
\begin{equation}
(W^E - W^D)[Z_\mathrm{out},Z_\mathrm{in}] \;\geq\; 0 \qquad \forall\,Z_\mathrm{out},Z_\mathrm{in}\in\mathcal{P},
\end{equation}
equivalently $W^E v \geq W^D v$ entrywise for all $v\in\R^4_{\geq 0}$.
\end{lemma}
\begin{proof}
By Definition~\ref{def:row_sum_operators},
\[
W^E - W^D = W^{4K} + (d-1)\,W^{U_1} - W^{A} - W^{U_2}.
\]
By \eqref{eq:W-bulk}, we have \begin{equation}\label{eq:W-bulk}
W^{P}_{Z_{\mathrm{out}}, Z_{\mathrm{in}}}
\;=\;
N_{Z_{\mathrm{out}}}(m)
\!\!\sum_{K \,\in\, \mathcal{K}(P)}\!\! \mu^{X}_{P, K}
\!\!\!\!\sum_{\substack{K_{1}, K_{2} \in \mathcal{P} \\ K_{1} \wedge K_{2} \,=\, K}}\!\!\!\!
N_{K_{1}}(m_{1})\, N_{K_{2}}(m_{2})\,
\mu^{g}_{\,Z_{\mathrm{in}} \wedge K_{1},\, K_{2} \wedge Z_{\mathrm{out}}}.
\end{equation}
the only $P$-dependent factor of $W^{P}_{Z_\mathrm{out},Z_\mathrm{in}}$ is $\mu^X_{P,K}$. 
Since $\mu^X_{P,K}=0$ for $K\notin\mathcal{K}(P)$, we can sum over all $\mathcal{P}$ getting:
\[
(W^E - W^D)[Z_\mathrm{out},Z_\mathrm{in}]
= N_{Z_\mathrm{out}}(m)\!\!\sum_{K\in\mathcal{P}}\Delta^X(K)\!\!\!\sum_{\substack{K_1,K_2\in\mathcal{P}\\K_1\wedge K_2=K}}\!\!\!N_{K_1}(m_1)N_{K_2}(m_2)\,\mu^g_{Z_\mathrm{in}\wedge K_1,\,K_2\wedge Z_\mathrm{out}},
\]
with
\[
\Delta^X(K):=\mu^X_{4K,K}+(d-1)\mu^X_{U_1,K}-\mu^X_{A,K}-\mu^X_{U_2,K}.
\]
We get the following
\begin{equation}\label{eq:T-diff}
(W^E - W^D)[Z_\mathrm{out},Z_\mathrm{in}]
= N_{Z_\mathrm{out}}(m)\!\!\sum_{K_1,K_2\in\mathcal{P}}\!\! N_{K_1}(m_1)N_{K_2}(m_2)\,\Delta^X(K_1\wedge K_2)\,\mu^g_{Z_\mathrm{in}\wedge K_1,\,K_2\wedge Z_\mathrm{out}}.
\end{equation}
Plugging in the moments of Lemma~\ref{lem:X-leaf-full},
\[
\Delta^X(K)=
\begin{cases}
\sigma_4 + (d-3)\sigma_2^2  & K=4K,\\[2pt]
d\,\sigma_2^2, & K=U_1,\\[2pt]
0, & K\in\{A,U_2\}.
\end{cases}
\]
By the assumptions of the lemma we get that $\Delta^X(K)\ge0$ for every $K\in \mathcal{P}$.
\end{proof}

\paragraph{Boundary lemmas.}
\begin{lemma}[End-vector dominance]\label{lem:end_lemma_dominance_max_row}
$u_\mathrm{end}^{E} - u_\mathrm{end}^{D} \geq 0$ entrywise.
\end{lemma}
\begin{proof}
Similar to the previous lemma. By \eqref{eq:u-end} the only $P_N$-dependent factor of $[u^{P_N}_\mathrm{end}]_{Z_\mathrm{in}}$ is $\mu^X_{P_N,\cdot}$. By Definition~\ref{def:row_sum_operators},
\[
u_\mathrm{end}^E-u_\mathrm{end}^D=u^{4K}_\mathrm{end}+(d-1)u^{U_1}_\mathrm{end}-u^A_\mathrm{end}-u^{U_2}_\mathrm{end},
\]
so combining the four copies of \eqref{eq:u-end},
\begin{equation}\label{eq:uend-diff}
[u_\mathrm{end}^E - u_\mathrm{end}^D]_{Z_\mathrm{in}}
= \!\!\sum_{\mathcal N\in\{A,4K\}}\!\! N_{\mathcal N}(m)\!\!\sum_{K_1,K_2\in\mathcal{P}}\!\! N_{K_1}(m_1)N_{K_2}(m_2)\,\Delta^X(K_1\wedge K_2)\,\mu^g_{Z_\mathrm{in}\wedge K_1,\,K_2\wedge\mathcal N}
\end{equation}
with $$\Delta^X(J):=\mu^X_{4K,J}+(d-1)\mu^X_{U_1,J}-\mu^X_{A,J}-\mu^X_{U_2,J}$$

Every factor is non-negative, and $\Delta^X(J)\geq 0$ as in the previous lemma, so $[u_\mathrm{end}^E - u_\mathrm{end}^D]_{Z_\mathrm{in}}\geq 0$.



\end{proof}

\begin{lemma}[Boundary dominance]\label{lem:uinit_max_row}
For $\sigma_4\ge\sigma_2^2$ and $d\ge2$,
\[
u_\mathrm{init}^{(E,E)} \;\ge\; u_\mathrm{init}^{(s_1,s_2)} \qquad\text{entrywise, for every }(s_1,s_2)\in\{E,D\}^2 .
\]
\end{lemma}
\begin{proof}
Fix $Z_\mathrm{out}\in\mathcal P$.
Denote $$G(K,J)\;:=\;N_K(m)\!\!\sum_{\substack{K_1,K_2\in\mathcal P\\ K_1\wedge K_2=J}}\!\!
N_{K_1}(m_1)\,N_{K_2}(m_2)\,\mu^g_{\,K\wedge K_1,\;K_2\wedge Z_\mathrm{out}}.$$
Substitute \eqref{eq:u-init} into definition~\ref{def:row_sum_operators} (both sums over all of $\mathcal P$, since
$\mu^X_{P,\cdot}=0$ off $\mathcal K(P)$) and subtract:
\begin{align}
[u_\mathrm{init}^{(E,E)}-u_\mathrm{init}^{(E,D)}]_{Z_\mathrm{out}}
&= N_{Z_\mathrm{out}}(m)\!\sum_{K,J\in\mathcal P}\! G(K,J)\,
\bigl(\mu^X_{4K,K}+(d-1)\mu^X_{U_1,K}\bigr)\nonumber\\
&\quad
\bigl[\bigl(\mu^X_{4K,J}+(d-1)\mu^X_{U_1,J}\bigr)-\bigl(\mu^X_{A,J}+\mu^X_{U_2,J}\bigr)\bigr],
\label{eq:dED}\\[4pt]
[u_\mathrm{init}^{(E,E)}-u_\mathrm{init}^{(D,E)}]_{Z_\mathrm{out}}
&= N_{Z_\mathrm{out}}(m)\!\sum_{K,J\in\mathcal P}\! G(K,J)\,
\bigl[\bigl(\mu^X_{4K,K}+(d-1)\mu^X_{U_1,K}\bigr)-\bigl(\mu^X_{A,K}+\mu^X_{U_2,K}\bigr)\bigr]\nonumber\\
&\quad
\bigl(\mu^X_{4K,J}+(d-1)\mu^X_{U_1,J}\bigr),
\label{eq:dDE}\\[4pt]
[u_\mathrm{init}^{(E,E)}-u_\mathrm{init}^{(D,D)}]_{Z_\mathrm{out}}
&= N_{Z_\mathrm{out}}(m)\!\sum_{K,J\in\mathcal P}\! G(K,J)\,
\Bigl[\bigl(\mu^X_{4K,K}+(d-1)\mu^X_{U_1,K}\bigr)\nonumber\\
&\quad
\bigl[\bigl(\mu^X_{4K,J}+(d-1)\mu^X_{U_1,J}\bigr)-\bigl(\mu^X_{A,J}+\mu^X_{U_2,J}\bigr)\bigr]\nonumber\\
&\quad
+\bigl[\bigl(\mu^X_{4K,K}+(d-1)\mu^X_{U_1,K}\bigr)-\bigl(\mu^X_{A,K}+\mu^X_{U_2,K}\bigr)\bigr]\nonumber\\
&\quad
\bigl(\mu^X_{A,J}+\mu^X_{U_2,J}\bigr)\Bigr].
\label{eq:dDD}
\end{align}

The factors $N_{Z_\mathrm{out}}(m),N_K(m),N_{K_1}(m_1),N_{K_2}(m_2)\ge0$, the gadget
moments $\mu^g_{\cdot,\cdot}\ge0$ (Lemma~\ref{lem:mu-g}), and the un-subtracted leaf
weights $\mu^X_{4K,\cdot}+(d-1)\mu^X_{U_1,\cdot}\ge0$ and
$\mu^X_{A,\cdot}+\mu^X_{U_2,\cdot}\ge0$. It remains to show the subtracted factor is
nonnegative, for every $B\in\mathcal P$:

By Lemma~\ref{lem:X-leaf-full}, with $\sigma_4+(d-3)\sigma_2^2\ge(d-2)\sigma_2^2\ge0$
for $\sigma_4\ge\sigma_2^2$, $d\ge2$.
\[
\bigl(\mu^X_{4K,B}+(d-1)\mu^X_{U_1,B}\bigr)-\bigl(\mu^X_{A,B}+\mu^X_{U_2,B}\bigr)
=
\begin{cases}
\sigma_4+(d-3)\sigma_2^2 & B=4K,\\
d\,\sigma_2^2            & B=U_1,\\
0                       & B\in\{A,U_2\},
\end{cases}
\qquad\ge\ 0,
\]
 Hence every summand of
\eqref{eq:dED}--\eqref{eq:dDD} is $\ge0$, so all three differences are $\ge0$. As
$Z_\mathrm{out}$ was arbitrary, $u_\mathrm{init}^{(E,E)}\ge u_\mathrm{init}^{(s_1,s_2)}$
entrywise for every $(s_1,s_2)\in\{E,D\}^2$.
\end{proof}

\begin{lemma}[Single-flip monotonicity]\label{lem:flip}
Let $N\ge3$, $\sigma_4\ge\sigma_2^2$, $d\ge2$ and let $\mathbf{s}\in\{E,D\}^N$.
From Theorem~\ref{thm:row-decomp},
\[
R(\mathbf{s})=\bigl\langle u_{end}^{(s_N)},\,W^{(s_{N-1})}\cdots W^{(s_3)}\,u_{init}^{(s_1,s_2)}\bigr\rangle
\]
is a product of $N-1$ elements: the boundary vector $u_{init}^{(s_1,s_2)}$, the $N-3$ transition matrices $W^{(s_3)},\dots,W^{(s_{N-1})}$, and the boundary vector $u_{end}^{(s_N)}$. If any one of these elements is of type $D$, replacing that element by its type-$E$ counterpart (and leaving all others unchanged) yields a sequence
$\mathbf{s}'$ with
\[
R(\mathbf{s}')\ \ge\ R(\mathbf{s}).
\]
\end{lemma}

\begin{lemma}[Single-flip monotonicity]\label{lem:flip}
Let $N\ge3$, $\sigma_4\ge\sigma_2^2$, $d\ge2$ and let $\mathbf{s}\in\{E,D\}^N$.
From theorem \ref{thm:row-decomp} we get a product of $N-1$  elements of the form $$R(\mathbf{s})=\langle u_{end}^{(\mathbf{s}^N)},W^{(\mathbf{s}^{N-1})}\cdots W^{(3)}u_{init}^{(\mathbf{s}_1,\mathbf{s}_2)}\rangle$$
For each element 
Let $N\ge3$, $\sigma_4\ge\sigma_2^2$, $d\ge2$ and let $\mathbf{s}\in\{E,D\}^N$.
From Theorem~\ref{thm:row-decomp} we get a product of $N-1$ elements of the form
\[
R(\mathbf{s})=\bigl\langle u_{end}^{(\mathbf{s}_N)},\,W^{(\mathbf{s}_{N-1})}\cdots W^{(\mathbf{s}_3)}\,u_{init}^{(\mathbf{s}_1,\mathbf{s}_2)}\bigr\rangle .
\]
If  $\mathbf{s}$ has an element of type $D$ let $\mathbf{s}'$ be the sequence obtained by replacing that element $D$ with the corresponding $E$, keeping all other elements fixed, we get $$R(\mathbf{s}^\prime)\ge R(\mathbf{s})$$
\end{lemma}
\begin{proof}

By \eqref{eq:row-transfer},
\[
R(\mathbf{s})=\bigl\langle u_\mathrm{end}^{(s_N)},\ W^{(s_{N-1})}\cdots W^{(s_3)}\,u_\mathrm{init}^{(s_1,s_2)}\bigr\rangle .
\]

Element in place $l=1,...,N$  in the sequence $s$  can be of type $D$  in 3 different places:
\smallskip\noindent\textbf{(i) Boundary $u_\mathrm{init}$ - $l\in\{1,2\}$.}
By Lemma~\ref{lem:uinit_max_row}, $u_\mathrm{init}^{(E,E)}\ge u_\mathrm{init}^{(s_1,s_2)}\ge0$. And we know that $ u_\mathrm{end}^{(s_N)},\ W^{(s_{N-1})}\cdots W^{(s_3)}$ all contain non negative element and hence by taking a sequence $\mathbf{s^\prime}=s^N\cdots s^3EE$, i.e. taking the element $u_{init}$  to be of type $u_{init}^{(E,E)}$  we get that
\[
R(\mathbf{s})\ \le\ R(\mathbf{s}^\prime).
\]

\smallskip\noindent\textbf{(ii) Interior $W$ - $l\in\{3,\dots,N-1\}$.}
By Lemma \ref{lem:transition_matrix_dominance}, $W^{E}\ge W^{(s_l)}\ge0$ entrywise. And we know that $u_\mathrm{end}^{(s_N)},\ W^{(s_{N-1})}\cdots W^{(s_{l+1})}$ and
$W^{(s_{l-1})}\cdots W^{(s_3)}\,u_\mathrm{init}^{(s_1,s_2)}$ all contain non negative elements, and hence by taking the sequence $\mathbf{s}^\prime$ that replaces the element
$W^{(s_l)}$ by one of type $W^{E}$ we get that
\[
R(\mathbf{s})\ \le\ R(\mathbf{s}^\prime).
\]

\smallskip\noindent\textbf{(iii) Boundary $u_\mathrm{end}$ - $l=N$.}
By Lemma~\ref{lem:end_lemma_dominance_max_row}, $u_\mathrm{end}^{E}\ge u_\mathrm{end}^{(s_N)}\ge0$.
And we know that $W^{(s_{N-1})}\cdots W^{(s_3)}\,u_\mathrm{init}^{(s_1,s_2)}$ all contain non negative elements, and hence by taking the sequence $\mathbf{s}^\prime$ that replaces the element $u_\mathrm{end}^{(s_N)}$ by one of type $u_\mathrm{end}^{E}$ we get that
\[
R(\mathbf{s})\ \le\ R(\mathbf{s}^\prime).
\]
We got that in all $3$  possible cases it is best to use an element of type $E$ .
\end{proof}
\begin{corollary}
\label{cor:max_row_sum}
A row with with E/D pattern $\mathbf{s}$ , see \ref{def:E_D_sites},  with $\mathbf{s}=E^N$ obtains the maximal row sum $R(E^N)$.
\end{corollary}
\begin{proof}
Assume by contradiction that this isn't true, so a sequence $\mathbf{s}^\prime$  with an element $D$  obtains the maximal sum and $$R(\mathbf{s}^\prime)>R(E^N)$$
But with the previous lemma we can inductively get that $R(\mathbf{s}^\prime)\le R(E^N)$ in contradiction to the assumption.
\end{proof}

\paragraph{Closed form of $R(E^N)$.}
By the previous corollary the maximum row sum equals $R(E^N)$, which we now express in closed form.
\begin{theorem}[Closed form of the maximal row sum]\label{thm:R-EN-closed}
Let $N\ge3$, $\sigma_4\ge\sigma_2^2$, $d\ge2$. Define the aggregated transfer matrix and
boundary vectors
\[
\mathbf{W}^{\mathrm{Max}} \;:=\; W^{4K} + (d-1)\,W^{U_1},
\qquad
\mathbf{u}^{\mathrm{Max}}_{\mathrm{end}} \;:=\; u_\mathrm{end}^{4K} + (d-1)\,u_\mathrm{end}^{U_1},
\]
\[
\mathbf{u}^{\mathrm{Max}}_{\mathrm{init}} \;:=\; u_\mathrm{init}^{(4K,4K)} + (d-1)\,u_\mathrm{init}^{(4K,U_1)}
+ (d-1)\,u_\mathrm{init}^{(U_1,4K)} + (d-1)^{2}\,u_\mathrm{init}^{(U_1,U_1)},
\]
where $W^{P},u_\mathrm{end}^{P},u_\mathrm{init}^{(P_1,P_2)}$ are given by
Theorems~\ref{thm:W-bulk},~\ref{thm:u-end},~\ref{thm:u-init}. Then the maximal row sum
admits the closed form
\[
R(E^N) \;=\; \bigl\langle \mathbf{u}^{\mathrm{Max}}_{\mathrm{end}},\; (\mathbf{W}^{\mathrm{Max}})^{\,N-3}\, \mathbf{u}^{\mathrm{Max}}_{\mathrm{init}}\bigr\rangle .
\]
\end{theorem}
\begin{proof}
From theorem \ref{thm:row-decomp}  and eq (\eqref{eq:row-decomp}) and using $V(\mathbf{P})$ from \eqref{eq:chain-value} 
\begin{align}
R(E^N)
&= \!\!\!\!\!\sum_{\mathbf{P}\in\{4K,U_1\}^N}\!\!\!\!\!(d-1)^{\#\{l:P_l=U_1\}}\, V(\mathbf{P}) \notag \\
&= \!\!\!\!\!\sum_{\mathbf{P}\in\{4K,U_1\}^N}\!\!\!\!\!(d-1)^{\#\{l:P_l=U_1\}}\, \bigl\langle u_\mathrm{end}^{P_N},\; W^{P_{N-1}}\!\cdots W^{P_3}\, u_\mathrm{init}^{(P_1,P_2)}\bigr\rangle. 
\end{align}
The multiplicity factor $(d-1)^{\#\{l:P_l=U_1\}} = \prod_l (d-1)^{[P_l=U_1]}$ splits across sites, and the inner product is multilinear in its $N$ arguments, so we may push each per-site sum inside:
\begin{equation}
R(E^N)
\;=\; \biggl\langle \!\!\sum_{P_N\in\{4K,U_1\}}\!\!\!(d-1)^{[P_N=U_1]} u_\mathrm{end}^{P_N}, \prod_{l=3}^{N-1}\biggl(\sum_{P_l\in\{4K,U_1\}}\!\!\!(d-1)^{[P_l=U_1]} W^{P_l}\biggr) \!\!\!\!\!\sum_{P_1, P_2\in\{4K,U_1\}}\!\!\!\!\!(d-1)^{[P_1=U_1]+[P_2=U_1]} u_\mathrm{init}^{(P_1,P_2)} \biggr\rangle.
\end{equation}
Defining the aggregated transfer matrix and boundary vectors
\begin{equation}
\mathbf{W}^{\mathrm{Max}} \;:=\; W^{4K} + (d-1)\,W^{U_1},
\qquad
\mathbf{u}^{\mathrm{Max}}_{\mathrm{end}} \;:=\; u_\mathrm{end}^{4K} + (d-1)\,u_\mathrm{end}^{U_1},
\end{equation}
\begin{equation}
\mathbf{u}^{\mathrm{Max}}_{\mathrm{init}} \;:=\; u_\mathrm{init}^{(4K,4K)} + (d-1)\,u_\mathrm{init}^{(4K,U_1)} + (d-1)\,u_\mathrm{init}^{(U_1,4K)} + (d-1)^{2}\,u_\mathrm{init}^{(U_1,U_1)},
\end{equation}
Plug that in the previous equation and get the required results
\begin{equation}
\;R(E^N) \;=\; \bigl\langle \mathbf{u}^{\mathrm{Max}}_{\mathrm{end}},\; (\mathbf{W}^{\mathrm{Max}})^{\,N-3}\, \mathbf{u}^{\mathrm{Max}}_{\mathrm{init}}\bigr\rangle.\;
\end{equation}

\end{proof}

\subsubsection{$R(E^N)$  is an eigenvalue}
\begin{lemma}\label{lem:E_N_EIGENVALUE}
Let $v\in\mathbb{R}^{d^{2N}}$ be the indicator of the column diagonal,
\[
v_{(i^{(b)},i^{(d)})} \;=\; \mathbb{1}\!\left[i^{(b)}=i^{(d)}\right]
.
\]
Then $v$ is an eigenvector of $M$ with eigenvalue $R(E^{N})$:
\[
M v \;=\; R(E^{N})\,v

\]
\end{lemma}

\begin{proof}
For any row $(i^{(a)},i^{(c)})$, using $v_{(i^{(b)},i^{(d)})}=\mathbb{1}[i^{(b)}=i^{(d)}]$
\begin{equation}\label{eq:Mv-is-Gsq}
(Mv)_{(i^{(a)},i^{(c)})}
=\sum_{i^{(b)}}M_{(i^{(a)},i^{(c)}),(i^{(b)},i^{(b)})}
\end{equation}
We sum over each row $(i^{(a)},i^{(c)})$, which leads to $2$ options:

 \begin{enumerate}
     \item $i^{(a)}\neq i^{(c)}$ - Because $i^{(a)}\not=i^{(c)}$  we get that when looking at Row classification by $E/D$ (\ref{def:E_D_sites}) we have some site $l=1,...,N$  with classification $D$. From lemma (\ref{lem:per_site_compa}) we know that each of these element are non zero only in columns where the pattern chain received has $P_l\in\{A,U_2\}$. But we sum only over columns where $i^{(b)}=i^{(d)}$  meaning $P_l\in\{4K,U_1\}$. So for this option the sum vanishes.
\item $i^{(a)}=i^{(c)}$ - From corollary (\ref{cor:EN-rowsum}) we get that the sum of the row equals $R(E^N)$
\end{enumerate}

Combining the two cases,
\[
(Mv)_{(i^{(a)},i^{(c)})}
\;=\;
\begin{cases}
R(E^N), & i^{(a)}=i^{(c)},\\[2pt]
0,      & i^{(a)}\neq i^{(c)},
\end{cases}
\;=\;
R(E^N)\,\mathbb{1}\!\left[i^{(a)}=i^{(c)}\right]
\;=\;
R(E^N)\,v_{(i^{(a)},i^{(c)})}.
\]
Hence $Mv = R(E^N)\,v$.
\end{proof}
We will now prove theorem \ref{thm:opnorm-closed}
\begin{proof}
Since $M$ is symmetric with nonnegative entries, Gershgorin's theorem (applied
rowwise) bounds every eigenvalue, $\lambda$, of $M$ by the maximal row sum,
$$\lambda \le \max_{\mathbf{s}\in\{E,D\}^N} R(\mathbf{s})$$
By corollary~\ref{cor:max_row_sum} this maximum equals $R(E^N)$.
Hence
$$\|M\|_\infty = \lambda_{max} \le R(E^N)$$
By Lemma~\ref{lem:E_N_EIGENVALUE}, $R(E^N)$ is itself an eigenvalue of $M$, so the bound
is attained and $\|M\|_\infty = R(E^N)$.
The closed form \eqref{eq:opnorm_formula} now follows from theorem~\ref{thm:R-EN-closed}.
\end{proof}


\end{document} 


\end{document}